\documentclass{article}
\usepackage[preprint]{neurips_2026}
\usepackage{amsmath,amssymb,amsfonts}
\usepackage{amsthm}
\usepackage{booktabs}
\usepackage{enumitem}
\usepackage{graphicx}
\usepackage[hypertexnames=false]{hyperref}
\usepackage{xcolor}
\usepackage{algorithm}
\usepackage{algpseudocode}
\setcitestyle{numbers}
\renewcommand{\theHtable}{tbl.\arabic{table}}

\newtheorem{theorem}{Theorem}[section]
\newtheorem{proposition}[theorem]{Proposition}
\newtheorem{corollary}[theorem]{Corollary}
\newtheorem{definition}[theorem]{Definition}
\newtheorem{remark}[theorem]{Remark}
% TODO tracking: search for \todo{} to find all placeholders
\newcommand{\todo}[1]{\textcolor{red}{[TODO: #1]}}

\title{Null-Space Filtered Low-Rank Adaptation\\
for CLIP Class-Incremental Learning}

\author{%
Xuan Rao$^1$, Yong Zeng$^2$, Bo Zhao$^1$, Derong Liu$^3$, and Cesare Alippi$^{4,5}$\\
$^1$School of Systems Science, Beijing Normal University, Beijing 100875, China\\
$^2$School of Automation and Intelligent Manufacturing, Southern University of Science and Technology, Shenzhen 518055, China\\
$^3$School of Artificial Intelligence, Anhui University, Hefei 230601, China\\
$^4$Dipartimento di Elettronica e Informazione, Politecnico di Milano, Milano 20133, Italy\\
$^5$Faculty of Informatics, Universit\`{a} della Svizzera Italiana, Lugano 6900, Switzerland\\
\texttt{raoxuan88@bnu.edu.cn}, \texttt{zhaobo@bnu.edu.cn}\\
\texttt{12532847@mail.sustech.edu.cn}, \texttt{derong@ahu.edu.cn}\\
\texttt{cesare.alippi@usi.ch}
}

\begin{document}

\maketitle

\begin{abstract}
Continually adapting CLIP to new classes can improve downstream discrimination while interfering with historical task directions and eroding the image--text relations acquired during pre-training. Standard low-rank adaptation limits the rank of an update but does not control the historical input directions along which adaptation occurs. We introduce \emph{LoRA-NF}, a low-rank adaptation method with persistent null-space filtering. LoRA-NF inserts a history-induced filter before both trainable low-rank factors, so that the adapter learns from attenuated historical high-energy directions throughout optimization. A leaky construction retains an adaptation path in every direction while biasing learning toward the approximate null space. We complement this historical-task protection with cross-modal knowledge distillation, which explicitly preserves image--text relations on reference pairs. At prediction time, a low-rank regularized Gaussian discriminant classifier is combined with the CLIP text classifier to integrate supervised visual statistics with language-derived semantic priors. We evaluate the resulting framework on class-incremental classification and image--text retrieval, separating the roles of forward filtering, cross-modal distillation, and ensemble prediction. \todo{Insert verified headline classification and retrieval results.}
\end{abstract}

%==============================================================================
\section{Introduction}
\label{sec:intro}

Vision-language models (VLMs), exemplified by CLIP~\cite{radford2021learning}, acquire broad visual concepts and open-vocabulary recognition capabilities by aligning images and natural language at scale. These properties make CLIP an appealing foundation for continual learning, where a shared model must absorb sequentially arriving classes without being retrained from scratch. Yet continual adaptation also changes what must be preserved. Beyond retaining discrimination over previously observed classes, an adapted CLIP model should maintain the image--text alignment that supports zero-shot recognition and cross-modal retrieval. We therefore study CLIP class-incremental learning through two complementary objectives: learning strong classifiers for sequentially introduced classes and preserving the cross-modal capability inherited from pre-training.

These objectives expose two related but distinct forms of degradation. Supervision from a new task reshapes the representation toward the currently observed classes and can interfere with directions that are important to historical tasks. Repeated adaptation can also move the representation away from the image--text relations learned during pre-training~\cite{zheng2023zscl,ni2023modx}. A frozen CLIP model retains its pre-trained geometry but cannot fully exploit new supervision, whereas unconstrained adaptation improves current-task discrimination at the cost of historical knowledge and cross-modal transfer. A continual learner must therefore control interference with historical task directions while explicitly anchoring cross-modal relations during training.

Low-rank adaptation (LoRA)~\cite{hu2022lora} provides a natural residual pathway for updating a frozen foundation model. For a linear transformation, LoRA parameterizes the residual update as $\Delta\mathbf{W}=\mathbf{A}\mathbf{B}$, with a rank much smaller than the ambient feature dimension. This factorization reduces the number of trainable parameters, but low rank alone is not a knowledge-preservation constraint: the input-side factor remains free to amplify directions on which historical tasks exhibit high activation energy. In other words, LoRA determines how many degrees of freedom are used for adaptation, but not which historical input directions these degrees of freedom should follow.

Null-space methods provide a complementary strategy by identifying directions that minimally affect previously acquired knowledge. Gradient-projection approaches modify an update after conventional forward and backward computation so that it lies in an approximate null space of historical features~\cite{saha2021gradient,kang2025dynamic}. Other null-space-inspired LoRA variants initialize or constrain the input-side factor using a basis derived from low-energy activation directions~\cite{tang2026lora}. These designs inject historical geometry into the gradient or a low-rank factor, but they do not place a fixed history-induced filter persistently before both trainable factors. This distinction motivates our central question: can historical activation geometry be embedded directly into the LoRA forward function, such that both low-rank factors are optimized in a filtered input space throughout adaptation?

We answer this question with \textbf{LoRA-NF}, low-rank adaptation with null-space filtering. Under the right-multiplication convention used throughout this paper, LoRA-NF replaces the conventional update $\mathbf{A}\mathbf{B}$ with
\begin{equation}
    \Delta\mathbf{W}=\mathbf{P}\mathbf{A}\mathbf{B},
    \label{eq:intro_lora_nf}
\end{equation}
where $\mathbf{P}$ is constructed from the eigenspectrum of historical layer-input second moments. The adapter consequently computes $(\mathbf{X}\mathbf{P})\mathbf{A}\mathbf{B}$: historical high-energy directions are attenuated before the two trainable factors learn their transformation. Unlike an initialization-only bias, the filter remains active throughout optimization. We use a leaky rather than strict filter, retaining an adaptation path in every direction while biasing learning away from directions most strongly occupied by historical tasks. The formal conditions governing its historical response, expressivity, and optimization dynamics are developed in Section~\ref{sec:lora_nf_theory}.

LoRA-NF is designed to preserve historical task knowledge; it is not, by itself, a mechanism for preserving cross-modal retrieval. We assign the latter role primarily to cross-modal knowledge distillation, which constrains the adapted model to reproduce the image--text relations of a frozen CLIP teacher on reference pairs. At inference, we further combine the CLIP text classifier with low-rank regularized Gaussian discriminant analysis (LR-RGDA), which captures class-conditional statistics in the adapted visual space. These branches provide complementary language-derived priors and supervised visual evidence. The prediction module contributes to continual classification, but does not affect image--text retrieval.

We evaluate the framework under the X-TAIL protocol~\cite{xu2024rail} across ten heterogeneous image-classification datasets and assess both class-incremental learning and image--text retrieval. The experiments examine whether LoRA-NF improves over rank- and budget-matched LoRA, projected-gradient, and null-space-inspired factor variants; whether cross-modal distillation preserves retrieval after sequential adaptation; whether LR-RGDA and the CLIP text classifier provide complementary decision signals; and whether the resulting classification gains coexist with preserved cross-modal capability. \todo{Insert verified Transfer/Average/Last, classifier-ablation, and retrieval results.}

Our contributions are summarized as follows:
\begin{itemize}[leftmargin=*]
    \item We introduce \textbf{LoRA-NF}, which reparameterizes the LoRA update from $\mathbf{A}\mathbf{B}$ to $\mathbf{P}\mathbf{A}\mathbf{B}$ and persistently filters adapter inputs using historical activation spectra. We characterize its historical response, expressive equivalence to rank-matched LoRA under an invertible leaky filter, and spectrum-dependent optimization geometry.
    \item We build a complete continual-learning framework around LoRA-NF, using cross-modal knowledge distillation to preserve image--text relations and an LR-RGDA/text ensemble to combine supervised visual statistics with CLIP's semantic prior.
    \item We systematically evaluate continual classification, cross-modal retrieval, mechanism-level predictions, training and classifier ablations, and the resulting adaptation--preservation behavior under a unified protocol.
\end{itemize}

\section{Related Work}
\label{sec:related}

\subsection{Class-Incremental Learning with Pre-trained Vision-Language Models}

Traditional class-incremental learning (CIL) requires a model to incorporate new classes while retaining discrimination over previously observed ones~\cite{de2021continual,wang2024b}. Representative strategies preserve knowledge through parameter regularization~\cite{kirkpatrick2017overcoming,ahn2019uncertainty}, knowledge distillation~\cite{li2017learning}, rehearsal or experience replay~\cite{rebuffi2017icarl,rolnick2019experience}, and parameter isolation or prompting~\cite{wang2022learning,wang2022dualprompt}. Prototype-based prediction and classifier correction further address the imbalance between old and newly introduced classes. These methods establish historical-task preservation as the central CIL objective, but generally assume that the representation itself has no independent pre-trained capability that must also be retained.

Pre-trained vision-language models extend this objective. In addition to old-class discrimination, continual adaptation should preserve open-vocabulary recognition and the image--text geometry acquired from large-scale pre-training. ZSCL~\cite{zheng2023zscl} constrains sequential adaptation to prevent zero-shot transfer degradation, while Mod-X~\cite{ni2023modx} analyzes intra-modal rotation and inter-modal deviation. C-CLIP~\cite{liu2025cclip} studies multimodal continual learning, and MG-CLIP~\cite{huang2025mgclip} explicitly preserves and compensates for the modality gap. RAIL~\cite{xu2024rail} and LADA~\cite{luo2025lada} further demonstrate the importance of prediction design under task-agnostic cross-domain evaluation. In particular, LADA freezes the visual encoder, adapts text representations, and combines a label-specific visual memory with a text classifier. Our framework instead continually adapts the representation and assigns historical-task preservation and cross-modal preservation to distinct training mechanisms.

\subsection{Knowledge-Preserving Low-Rank Adaptation}

Parameter-efficient fine-tuning adapts pre-trained models through a small set of task-dependent parameters. LoRA~\cite{hu2022lora} constrains residual weight updates to a low-rank factorization, while adapters, prompts~\cite{wang2022sprompts,yu2024moe}, and dynamic-rank variants~\cite{wang2024lora,lu2026codyra} control different aspects of adaptation capacity. Parameter efficiency, however, does not itself specify which historical representation directions may be modified. A low-rank update can still align with high-energy historical activations and substantially change the corresponding responses.

Subspace-based continual-learning methods instead protect historical knowledge by constraining update directions. GPM~\cite{saha2021gradient} projects gradients onto an approximate null space of previous-task features, and DMNSP~\cite{kang2025dynamic} extends multi-layer null-space projection to vision-language continual learning. Null-space-inspired LoRA methods also use low-energy activation directions to initialize or constrain the input-side factor~\cite{tang2026lora}. These approaches act on gradients or on the parameterization of an individual factor. LoRA-NF differs by placing a history-induced filter before both trainable low-rank factors, thereby maintaining the filtered forward form $(\mathbf{X}\mathbf{P})\mathbf{A}\mathbf{B}$ throughout optimization. Cross-modal distillation complements this directional protection by explicitly constraining image--text relations rather than historical activations alone.

\subsection{Continual Prediction with Pre-trained Semantic Priors}

Prototype, nearest-class-mean, and Gaussian discriminant classifiers summarize observed classes without repeatedly optimizing a parametric output layer~\cite{rebuffi2017icarl}. In pre-trained-model CIL, analytic classifiers and label memories can exploit stable feature statistics, whereas the CLIP text classifier provides a language-derived prior over both observed and unobserved classes. RAIL~\cite{xu2024rail} relies on the pre-trained text classifier for transfer, and LADA~\cite{luo2025lada} combines text predictions with a supervised label-memory branch. Gaussian discriminant analysis has also been applied to deep features~\cite{lee2018simple}, but class-specific full covariances are costly and poorly estimated in few-shot regimes. We use LR-RGDA to model regularized class-conditional structure with low-rank covariance corrections, and combine its supervised scores with the CLIP text classifier. This prediction module is evaluated separately from the representation-learning mechanisms.

\section{Method}
\label{sec:method}

\subsection{Problem Formulation and Framework Overview}
\label{sec:problem}

We consider a sequence of $T$ class-incremental tasks $\{\mathcal{D}_t\}_{t=1}^{T}$. Task $t$ contains labeled samples from a class set $\mathcal{C}_t$, where $\mathcal{C}_t\cap\mathcal{C}_{t'}=\varnothing$ for $t\neq t'$. At training step $t$, the learner receives $\mathcal{D}_t$ but no raw samples from earlier tasks. At inference, it predicts over the accumulated label space $\mathcal{C}_{1:t}=\bigcup_{j=1}^{t}\mathcal{C}_j$ without task identity. We use a pre-trained CLIP model with image encoder $f_I$ and text encoder $f_T$. The framework is evaluated on both class-incremental classification and image--text retrieval; the latter is computed directly from the adapted encoders and is independent of the downstream ensemble classifier.

Our framework separates three responsibilities. First, LoRA-NF uses historical layer-input statistics to reduce interference with directions important to earlier tasks. Second, cross-modal knowledge distillation explicitly anchors the image--text relations of the pre-trained model. Third, LR-RGDA and the CLIP text classifier combine supervised visual statistics with language-derived semantic priors at prediction time. After each task, the low-rank weights are merged into the backbone, the adapter factors are reinitialized, and the historical activation and class statistics are updated.

\paragraph{Notation.}
We adopt a row-vector, right-multiplication convention. For a linear layer,
\begin{equation}
    \mathbf{Y}=\mathbf{X}\mathbf{W},
    \qquad
    \mathbf{X}\in\mathbb{R}^{n\times d_{\mathrm{in}}},
    \quad
    \mathbf{W}\in\mathbb{R}^{d_{\mathrm{in}}\times d_{\mathrm{out}}}.
    \label{eq:linear_right}
\end{equation}
The implementation uses the transposed PyTorch storage convention, but all derivations below follow Eq.~\eqref{eq:linear_right}.

\subsection{LoRA-NF: Persistent Null-Space Filtering}
\label{sec:lora_nf}

Standard LoRA~\cite{hu2022lora} augments the frozen weight $\mathbf{W}_0$ by
\begin{equation}
    \mathbf{W}=\mathbf{W}_0+\Delta\mathbf{W},
    \qquad
    \Delta\mathbf{W}=\mathbf{A}\mathbf{B},
    \label{eq:lora}
\end{equation}
where $\mathbf{A}\in\mathbb{R}^{d_{\mathrm{in}}\times r}$, $\mathbf{B}\in\mathbb{R}^{r\times d_{\mathrm{out}}}$, and $r\ll\min(d_{\mathrm{in}},d_{\mathrm{out}})$. The residual response is $\mathbf{X}\mathbf{A}\mathbf{B}$. Although its rank is bounded by $r$, the factor $\mathbf{A}$ may extract arbitrary input directions, including directions that carry high energy for historical tasks.

\subsubsection{Historical activation spectrum}

For each adapted layer, let $\mathbf{X}_t\in\mathbb{R}^{N_t\times d_{\mathrm{in}}}$ collect its flattened input activations on task $t$. We compute the uncentered second moment
\begin{equation}
    \mathbf{G}_t=\frac{1}{N_t}\mathbf{X}_t^\top\mathbf{X}_t,
    \label{eq:task_gram}
\end{equation}
and maintain an equally weighted history over completed tasks,
\begin{equation}
    \overline{\mathbf{G}}_t=\frac{1}{t}\sum_{j=1}^{t}\mathbf{G}_j.
    \label{eq:history_gram}
\end{equation}
This matrix is stored per layer; raw historical inputs are not retained. Let
\begin{equation}
    \overline{\mathbf{G}}_t
    =\mathbf{U}\boldsymbol{\Lambda}\mathbf{U}^\top,
    \qquad
    0\leq\lambda_1\leq\cdots\leq\lambda_{d_{\mathrm{in}}},
    \label{eq:history_eig}
\end{equation}
be its eigendecomposition. The eigenvectors associated with small eigenvalues span a low-energy subspace in which historical activations have limited response. We collect the selected eigenvectors in $\mathbf{U}_0$ and define the leaky filter
\begin{equation}
    \mathbf{P}
    =(1-\rho)\mathbf{U}_0\mathbf{U}_0^\top+\rho\mathbf{I},
    \qquad 0<\rho\leq 1.
    \label{eq:leaky_filter}
\end{equation}
Thus $\mathbf{P}$ has eigenvalue $1$ on the selected low-energy subspace and $\rho$ on its complement. The first task uses $\mathbf{P}=\mathbf{I}$ because no historical activation spectrum is yet available. Soft spectral weights and strict projection are treated as variants in the ablation study; Eq.~\eqref{eq:leaky_filter} is the default LoRA-NF construction.

\subsubsection{Forward-filtered low-rank adaptation}

LoRA-NF inserts $\mathbf{P}$ before both trainable factors:
\begin{equation}
    \Delta\mathbf{W}=\mathbf{P}\mathbf{A}\mathbf{B},
    \qquad
    \mathbf{X}\Delta\mathbf{W}=(\mathbf{X}\mathbf{P})\mathbf{A}\mathbf{B}.
    \label{eq:lora_nf}
\end{equation}
The historical filter is fixed within a task, whereas $\mathbf{A}$ and $\mathbf{B}$ remain trainable. Consequently, the low-rank factors are optimized using a filtered input in every forward pass. This differs from modifying a gradient after ordinary forward/backward computation and from encoding a null-space basis only in the initialization or constraint of $\mathbf{A}$. After learning task $t$, we merge $\mathbf{P}\mathbf{A}\mathbf{B}$ into $\mathbf{W}_0$, reinitialize $\mathbf{A}$ and $\mathbf{B}$, collect $\mathbf{G}_t$, and construct the filter for the next task. This merge-and-reset procedure prevents trainable parameters from growing with the task sequence.

\subsection{Theoretical Understanding of LoRA-NF}
\label{sec:lora_nf_theory}

We next formalize three properties of the forward-filtered parameterization. The results describe a single adapted layer under the right-multiplication convention; complete proofs are provided in Appendix~\ref{app:lora_nf_proofs}.

\begin{proposition}[Historical filtered response]
\label{prop:historical_response}
For a historical activation matrix $\mathbf{X}_h$ with second moment $\mathbf{G}_h=\mathbf{X}_h^\top\mathbf{X}_h/N_h$, the filtered input satisfies
\begin{equation}
    \frac{1}{N_h}\|\mathbf{X}_h\mathbf{P}\|_F^2
    =\operatorname{tr}(\mathbf{P}^\top\mathbf{G}_h\mathbf{P}).
    \label{eq:filtered_energy}
\end{equation}
If $\mathbf{P}$ and $\mathbf{G}_h$ share eigenvectors with spectral weights $p_i$, then the right-hand side equals $\sum_i p_i^2\lambda_i$. Moreover,
\begin{equation}
    \|\mathbf{X}_h\mathbf{P}\mathbf{A}\mathbf{B}\|_F
    \leq \|\mathbf{X}_h\mathbf{P}\|_F\,\|\mathbf{A}\|_2\,\|\mathbf{B}\|_2.
    \label{eq:historical_bound}
\end{equation}
\end{proposition}

Proposition~\ref{prop:historical_response} connects the spectral filter to the response of the adapter on historical inputs. It is an interference bound rather than a guarantee of zero forgetting: the realized response also depends on the learned low-rank factors and on nonlinear interactions across layers.

\begin{proposition}[Expressive equivalence under leaky filtering]
\label{prop:expressive_equivalence}
Let $\mathbf{P}$ be invertible. For a fixed rank budget $r$,
\begin{equation}
    \{\mathbf{P}\mathbf{A}\mathbf{B}:\mathbf{A}\in\mathbb{R}^{d_{\mathrm{in}}\times r},
      \mathbf{B}\in\mathbb{R}^{r\times d_{\mathrm{out}}}\}
    =
    \{\mathbf{C}\mathbf{B}:\mathbf{C}\in\mathbb{R}^{d_{\mathrm{in}}\times r},
      \mathbf{B}\in\mathbb{R}^{r\times d_{\mathrm{out}}}\}.
    \label{eq:expressive_set}
\end{equation}
\end{proposition}

Because $\rho>0$ makes Eq.~\eqref{eq:leaky_filter} invertible, LoRA-NF does not obtain historical protection by shrinking the function class of rank-matched LoRA. The strict projection limit $\rho=0$ is excluded from this result.

\begin{proposition}[Spectrum-dependent optimization geometry]
\label{prop:optimization_geometry}
Define the effective input-side factor $\mathbf{C}=\mathbf{P}\mathbf{A}$ and suppose that $\mathbf{A}$ is updated by gradient descent with step size $\eta$ while $\mathbf{P}$ is fixed. Then
\begin{equation}
    \mathbf{C}^{+}
    =\mathbf{C}-\eta\mathbf{P}\mathbf{P}^\top\nabla_{\mathbf{C}}\mathcal{L}.
    \label{eq:optimization_geometry}
\end{equation}
For symmetric $\mathbf{P}$, the induced preconditioner is $\mathbf{P}^2$.
\end{proposition}

Equation~\eqref{eq:optimization_geometry} shows that the persistent filter changes the speed of effective learning along different historical eigendirections even though the expressive set is unchanged. For the leaky filter, directions outside the selected low-energy subspace receive a factor $\rho^2$. The statement is exact for gradient descent on $\mathbf{A}$; adaptive-optimizer dynamics are discussed separately and are not claimed to be identical.

\subsection{Cross-Modal Knowledge Distillation}
\label{sec:distillation}

LoRA-NF controls interference measured by historical activations, but it does not explicitly preserve CLIP's image--text alignment. We therefore use a frozen pre-trained CLIP model as teacher on a reference set of image--text pairs. Let $\mathbf{v}_i^T,\mathbf{t}_i^T$ denote normalized teacher features and $\mathbf{v}_i^S,\mathbf{t}_i^S$ the corresponding student features. Feature distillation anchors the visual representation through
\begin{equation}
    \mathcal{L}_{\mathrm{FD}}
    =\frac{1}{N}\sum_{i=1}^{N}\left(1-(\mathbf{v}_i^T)^\top\mathbf{v}_i^S\right).
    \label{eq:fd}
\end{equation}
To preserve relational image--text structure, we form temperature-scaled teacher and student distributions
\begin{equation}
    q_{ij}^{T}=\operatorname{softmax}_{j}\!\left(
      \frac{\exp(\gamma)(\mathbf{v}_i^T)^\top\mathbf{t}_j^T}{\tau}\right),
    \qquad
    q_{ij}^{S}=\operatorname{softmax}_{j}\!\left(
      \frac{\exp(\gamma)(\mathbf{v}_i^S)^\top\mathbf{t}_j^S}{\tau}\right),
    \label{eq:cross_modal_probs}
\end{equation}
where $\gamma$ is CLIP's logit-scale parameter. The default cross-modal distillation loss is the forward KL divergence
\begin{equation}
    \mathcal{L}_{\mathrm{CD}}
    =\frac{\tau^2}{N}\sum_{i=1}^{N}
      D_{\mathrm{KL}}(\mathbf{q}_i^T\,\|\,\mathbf{q}_i^S).
    \label{eq:cd}
\end{equation}
The complete training objective is
\begin{equation}
    \mathcal{L}
    =\mathcal{L}_{\mathrm{CE}}
    +\lambda_{\mathrm{FD}}\mathcal{L}_{\mathrm{FD}}
    +\lambda_{\mathrm{CD}}\mathcal{L}_{\mathrm{CD}}.
    \label{eq:training_objective}
\end{equation}
This separation is important: LoRA-NF primarily protects historical task directions, whereas cross-modal distillation is the mechanism directly responsible for maintaining image--text relations and is evaluated through retrieval.

\subsection{Distributional Ensemble Prediction}
\label{sec:ensemble}

After task $t$, we summarize each observed class $c\in\mathcal{C}_{1:t}$ by its feature mean $\boldsymbol{\mu}_c$ and covariance $\mathbf{\Sigma}_c$ in the adapted visual space. Few-shot class covariances are regularized as
\begin{equation}
    \mathbf{\Sigma}_c^{\mathrm{reg}}
    =\alpha_1\mathbf{\Sigma}_c
    +\alpha_2\mathbf{\Sigma}_{\mathrm{avg}}
    +\alpha_3\mathbf{I},
    \label{eq:cov_reg}
\end{equation}
where $\mathbf{\Sigma}_{\mathrm{avg}}$ is the dataset-balanced global covariance. The Gaussian discriminant score is
\begin{equation}
    g_c^{\mathrm{RGDA}}(\mathbf{z})
    =-\frac{1}{2}(\mathbf{z}-\boldsymbol{\mu}_c)^\top
      (\mathbf{\Sigma}_c^{\mathrm{reg}})^{-1}
      (\mathbf{z}-\boldsymbol{\mu}_c)
      -\frac{1}{2}\log|\mathbf{\Sigma}_c^{\mathrm{reg}}|+\log\pi_c.
    \label{eq:rgda}
\end{equation}

To retain class-specific structure without full quadratic cost, LR-RGDA represents the regularized covariance as
\begin{equation}
    \mathbf{\Sigma}_c^{\mathrm{reg}}
    =\mathbf{B}+\mathbf{U}_c\boldsymbol{\Lambda}_c\mathbf{U}_c^\top,
    \label{eq:lr_cov}
\end{equation}
where $\mathbf{B}=\alpha_2\mathbf{\Sigma}_{\mathrm{avg}}+\alpha_3\mathbf{I}$ is shared and $\mathbf{U}_c\boldsymbol{\Lambda}_c\mathbf{U}_c^\top$ is a rank-$r_g$ class-specific correction. Applying the Woodbury identity yields
\begin{equation}
    (\mathbf{\Sigma}_c^{\mathrm{reg}})^{-1}
    =\mathbf{B}^{-1}-\mathbf{B}^{-1}\mathbf{U}_c
      (\boldsymbol{\Lambda}_c^{-1}+\mathbf{U}_c^\top\mathbf{B}^{-1}\mathbf{U}_c)^{-1}
      \mathbf{U}_c^\top\mathbf{B}^{-1},
    \label{eq:woodbury}
\end{equation}
which decomposes the discriminant into a shared affine term and a class-specific low-rank quadratic correction. The full derivation and complexity analysis are provided in Appendices~\ref{app:lr_rgda_derivation} and~\ref{app:complexity}.

For completeness, define
\begin{equation}
    \mathbf{M}_c
    =\boldsymbol{\Lambda}_c^{-1}
    +\mathbf{U}_c^\top\mathbf{B}^{-1}\mathbf{U}_c.
\end{equation}
After dropping the class-independent quadratic term, the LR-RGDA score can be written as
\begin{equation}
    g_c^{\mathrm{LR\text{-}RGDA}}(\mathbf{z})
    =\mathbf{w}_c^\top\mathbf{z}+b_c+\mathbf{z}^\top\mathbf{Q}_c\mathbf{z},
    \label{eq:lr_rgda}
\end{equation}
where
\begin{align}
    \mathbf{w}_c
    &=\mathbf{B}^{-1}\boldsymbol{\mu}_c
      -\mathbf{B}^{-1}\mathbf{U}_c\mathbf{M}_c^{-1}
       \mathbf{U}_c^\top\mathbf{B}^{-1}\boldsymbol{\mu}_c,
       \label{eq:wc}\\
    b_c
    &=-\frac{1}{2}\boldsymbol{\mu}_c^\top\mathbf{B}^{-1}\boldsymbol{\mu}_c
      +\frac{1}{2}\boldsymbol{\mu}_c^\top\mathbf{B}^{-1}\mathbf{U}_c
       \mathbf{M}_c^{-1}\mathbf{U}_c^\top\mathbf{B}^{-1}\boldsymbol{\mu}_c
      -\frac{1}{2}\log|\mathbf{\Sigma}_c^{\mathrm{reg}}|+\log\pi_c,
      \label{eq:bc}\\
    \mathbf{Q}_c
    &=\frac{1}{2}\mathbf{B}^{-1}\mathbf{U}_c\mathbf{M}_c^{-1}
      \mathbf{U}_c^\top\mathbf{B}^{-1}.
      \label{eq:qc}
\end{align}

Let $s_c^{\mathrm{text}}(\mathbf{x})$ be the CLIP text-classifier logit and $s_c^{\mathrm{GDA}}(\mathbf{x})$ the LR-RGDA logit for an observed class. We first max-shift the two branches independently, producing $\widetilde{s}^{\mathrm{text}}$ and $\widetilde{s}^{\mathrm{GDA}}$. The ensemble uses
\begin{equation}
    s_c^{\mathrm{ens}}(\mathbf{x})=
    \begin{cases}
      (1-\alpha)\widetilde{s}_c^{\mathrm{text}}(\mathbf{x})
      +\alpha\widetilde{s}_c^{\mathrm{GDA}}(\mathbf{x}),
      & c\in\mathcal{C}_{1:t},\\
      (1-\alpha)\widetilde{s}_c^{\mathrm{text}}(\mathbf{x}),
      & c\notin\mathcal{C}_{1:t}.
    \end{cases}
    \label{eq:ensemble}
\end{equation}
The text branch can use cached adapted prototypes for observed classes and frozen CLIP prototypes for unobserved classes; we refer to it generically as the \emph{CLIP text classifier}. LR-RGDA changes only the classification decision and has no role in image--text retrieval.

\iffalse
\section{Introduction}
\label{sec:intro}
%------------------------------------------------------------------------------
% Paragraph 1: Task and Applications
Vision-language models (VLMs), exemplified by CLIP~\cite{radford2021learning}, have demonstrated remarkable zero-shot transfer capabilities across diverse downstream tasks by pre-training on large-scale image-text pairs. These models naturally serve as continual learners due to their ability to transfer representations across different domains. As real-world applications increasingly demand models that can sequentially learn from new tasks while retaining previously acquired knowledge, continual learning with CLIP has emerged as a critical research direction~\cite{wang2024lora,zheng2023zscl}.

%------------------------------------------------------------------------------
% Paragraph 2: Technical Challenge
However, adapting CLIP representations using only task-specific data can severely impair the general knowledge encoded in the pre-trained parameters~\cite{zheng2023zscl}. This challenge is particularly problematic in continual learning settings, where the model must not only learn incrementally from a series of tasks but also retain its pre-trained capabilities. Existing CLIP-based methods adapt the pre-trained image encoder by adding learnable parameters, but they often suffer from catastrophic forgetting~\cite{french1999catastrophic}---a phenomenon where the performance on previously learned tasks degrades as the model adapts to new ones. Furthermore, the trade-off between memory stability and learning plasticity~\cite{mermillod2013stability} becomes more complex when dealing with vision-language models due to their dual-modal nature. Beyond training-phase challenges, the inference-stage classifier design also plays a critical role: frozen zero-shot classifiers excel at generalization but underperform on in-distribution (ID) data, while supervised classifiers excel on seen classes but lack generalization to unseen domains.

%------------------------------------------------------------------------------
% Paragraph 3: Our Method
To address these limitations, we study a two-stage framework that separates representation adaptation from classifier construction in CLIP continual learning.

\textbf{Training Stage: LoRA-NSP.} At the training stage, we introduce \textbf{LoRA-NSP} (Low-Rank Adaptation with Null Space Parameterization), which incorporates soft-projected null space constraints into the low-rank adaptation process. Unlike conventional NSP that zeroes out signal subspaces, our soft projection preserves all eigendirections with adaptive weighting, achieving better stability-plasticity trade-offs. Additionally, we design a distillation loss combining feature distillation (FD) and cross-modal distillation (CD) to maintain representation stability.

\textbf{Inference Stage: Ensemble Classifier.} At the inference stage, we study an \textbf{ensemble classifier framework} that integrates LR-RGDA (Low-Rank Regularized Gaussian Discriminant Analysis) with the zero-shot classifier. LR-RGDA uses supervised class statistics extracted from the adapted CLIP feature space, while the zero-shot branch preserves CLIP's text-aligned classifier. This design is particularly relevant when previous task features cannot be stored and must be approximated by compact statistical replay.

%------------------------------------------------------------------------------
% Paragraph 4: Contributions
Our contributions are summarized as follows:
\begin{itemize}
    \item We establish a theoretical connection between null space projection (NSP) and low-rank adaptation (LoRA), showing that NSP can be viewed as a restricted form of LoRA with limited plasticity. Building on this insight, we propose LoRA-NSP with soft forward-pass projection for better stability-plasticity trade-offs, and design an FD+CD distillation mechanism to maintain representation stability and vision-language alignment during continual learning.
    \item We develop LR-RGDA, a regularized Gaussian discriminant classifier that exploits low-rank covariance structures via the Woodbury matrix identity, reducing inference complexity relative to full covariance discriminant analysis while retaining class-specific covariance corrections.
    \item We introduce an inference-time ensemble that combines LR-RGDA with zero-shot classification, and we evaluate its behavior under real-feature and lightweight statistical replay protocols against LADA-style label-memory classifiers.
\end{itemize}

%------------------------------------------------------------------------------
% Paragraph 5: Experiments
We evaluate our method on the recently proposed X-TAIL benchmark~\cite{xu2024rail}, which consists of 10 image classification datasets from diverse domains. Current verified evidence focuses on the inference-stage classifier comparison under matched feature sources and storage protocols; strict incremental training-stage results for LoRA-NSP are being evaluated separately and will be reported only after the same leakage checks and three-seed protocol are satisfied.

%==============================================================================
\section{Related Work}
\label{sec:related}

\textbf{Continual Learning Settings.}
Early continual learning research considers Task-Incremental Learning (TIL)~\cite{hsu2018re} and Class-Incremental Learning (CIL)~\cite{rebuffi2017icarl}, both assuming tasks from a single domain. Multi-domain Task-Incremental Learning (MTIL)~\cite{zheng2023zscl} relaxes this by introducing diverse domain-specific tasks. Cross-domain Task-Agnostic Incremental Learning (X-TAIL)~\cite{xu2024rail} removes task or domain identities entirely, representing the most challenging and practical setting. We conduct experiments under the X-TAIL protocol.

\textbf{Vision-Language Model Continual Learning.}
Existing CL methods for VLMs fall into three categories. \textit{Regularization-based} methods penalize parameter shifts to preserve pre-trained knowledge: ZSCL~\cite{zheng2023zscl} constrains the model in parameter space to prevent zero-shot transfer degradation. \textit{Architecture-based} methods introduce task-specific parameters: MoE-Adapters~\cite{yu2024moe} employ mixture-of-experts adapters with task-dependent routing, while LADA~\cite{luo2025lada} appends label-specific memory units to the frozen encoder. C-CLIP~\cite{liu2025cclip} establishes a multimodal benchmark and proposes an end-to-end framework balancing old and new task learning. \textit{Geometric analysis} approaches study representation structure: Mod-X~\cite{ni2023modx} identifies Spatial Disorder caused by intra-modal rotation and inter-modal deviation, proposing off-diagonal information preservation. MG-CLIP~\cite{huang2025mgclip} reveals the modality gap as a key factor and proposes gap preservation and compensation. Recently, \textit{parameter-efficient fine-tuning} has been explored for VLM continual learning: LoRA-CL~\cite{wang2024lora} studies low-rank adaptation in CL scenarios, and CoDyRA~\cite{lu2026codyra} dynamically optimizes LoRA ranks per parameter and task via sparsity-promoting regularization. Despite these advances, existing methods lack explicit theoretical grounding connecting parameter-efficiency constraints with knowledge preservation mechanisms.

\textbf{Null Space Projection for Catastrophic Forgetting.}
Null space projection (NSP) mitigates forgetting by constraining parameter updates to directions minimally affecting historical representations. GPM~\cite{saha2021gradient} projects gradients onto the approximate null space of previous task features, demonstrating strong stability at the cost of reduced plasticity. DMNSP~\cite{kang2025dynamic} extends this to vision-language continual learning by applying multi-layer null space projection exclusively to the visual branch with a dynamic projection coefficient, achieving finer stability-plasticity control. Concurrently, LoRA-Null~\cite{tang2026lora} initializes LoRA adapters in the null space of input activations to preserve pre-trained knowledge. However, these approaches treat NSP and LoRA as independent mechanisms: GPM and DMNSP perform hard gradient projection without leveraging LoRA's parameter-efficient structure, while LoRA-Null only affects initialization. Furthermore, all existing NSP methods employ hard projection that zeroes out signal subspaces, inherently limiting model plasticity.

\textbf{Inference-Time Classifier Design.}
Beyond training-phase adaptation, inference-time strategies are critical for balancing seen-class discrimination and zero-shot transfer. RAIL~\cite{xu2024rail} extends the classifier dimension while freezing feature representations, relying on the vanilla zero-shot CLIP classifier to handle unseen classes. Gaussian discriminant analysis on deep features has been explored for uncertainty estimation and OOD detection~\cite{lee2018simple}, but its direct application as a memory-efficient classifier in CLIP continual learning remains underexplored due to computational and sample-efficiency challenges. Our approach differs by integrating a supervised LR-RGDA classifier with the zero-shot classifier and by explicitly studying real-feature replay versus compact statistical replay.

\textbf{Relation to Our Work.}
Unlike existing methods that focus on either training or inference optimization, our approach studies how the two stages interact. \textit{Compared to NSP-based methods} (GPM, DMNSP), LoRA-NSP establishes a theoretical connection showing that NSP is a restricted form of LoRA, and proposes soft projection that preserves all eigendirections rather than hard zeroing. \textit{Compared to LoRA-Null}, LoRA-NSP operates throughout training via forward-pass projection rather than initialization only, and is validated on VLM continual learning. \textit{Compared to CoDyRA}, LoRA-NSP controls update directions via projection matrices rather than varying rank sizes, offering orthogonal and complementary benefits. \textit{Compared to RAIL and LADA}, our inference study emphasizes classifier construction from compact statistics, which targets settings where storing historical real features or label-specific memory is undesirable. The training innovations (LoRA-NSP + FD + CD) and inference innovations (LR-RGDA ensemble) are orthogonal and can be combined, but we evaluate their contributions separately to avoid conflating representation gains with classifier gains.

%==============================================================================
\section{Method}
\label{sec:method}

\subsection{Preliminaries}

\textbf{Problem Setting.}
Cross-domain Task-Agnostic Incremental Learning (X-TAIL)~\cite{xu2024rail} is defined as the scenario where given a pre-trained vision-language model, the learner is required to incrementally learn $K$ different tasks $\{\mathcal{T}^1, \mathcal{T}^2, ..., \mathcal{T}^K\}$. Each task $\mathcal{T}^k = (\mathcal{D}^k, C^k)$ is available only during the $k$-th learning step, where $\mathcal{D}^k = \bigcup_{j=1}^{M^k} \mathcal{D}_j^k$ with $M^k$ denotes the number of classes for task $\mathcal{T}^k$. During inference at all steps, the learner attempts to classify input images from any task without the task-identity hint.

\textbf{CLIP Model.}
The CLIP~\cite{radford2021learning} model contains an image encoder $f_I$ and a text encoder $f_T$. The training of CLIP is based on contrastive learning, which aligns image and text features in a shared latent space. For image classification, each class $j$ is transformed into a sentence $t_j^i$ using a template ``a photo of a $\{c_j^i\}$''. The text encoder $f_T$ processes $t_j^i$ to a text feature, while the image encoder $f_I$ processes the input image $v$, generating an image embedding. The class with the largest similarity score between the image embedding and text embeddings is the prediction for the image.

\subsection{Training Stage: LoRA-NSP}

\textbf{Motivation.}
Standard parameter-efficient fine-tuning methods for CLIP, such as LoRA~\cite{hu2022lora}, introduce low-rank decomposition matrices to adapt pre-trained models. However, these methods lack explicit constraints to protect pre-trained knowledge during continual learning, leading to catastrophic forgetting. To address this, we propose LoRA-NSP (Low-Rank Adaptation with Null Space Parameterization), which integrates null space constraints into the low-rank adaptation framework. Our approach is theoretically grounded on the insight that NSP can be viewed as a restricted form of LoRA, enabling us to design a soft projection mechanism that generalizes both approaches.

\textbf{Preliminaries: LoRA Parameterization.}
Given a pre-trained weight matrix $\mathbf{W}_0 \in \mathbb{R}^{d_{\text{out}} \times d_{\text{in}}}$, standard LoRA reparameterizes the weight update as:
\begin{equation}
    \mathbf{W} = \mathbf{W}_0 + s\mathbf{B}\mathbf{A}
    \label{eq:lora_param}
\end{equation}
where $\mathbf{B} \in \mathbb{R}^{d_{\text{out}} \times r}$, $\mathbf{A} \in \mathbb{R}^{r \times d_{\text{in}}}$ are trainable low-rank matrices with rank $r \ll \min(d_{\text{out}}, d_{\text{in}})$, and $s$ is a learnable or fixed scaling factor. By freezing $\mathbf{W}_0$ and only optimizing $\mathbf{A}$ and $\mathbf{B}$, LoRA drastically reduces the number of trainable parameters.

Under the LoRA parameterization, the gradients with respect to the trainable matrices are:
\begin{equation}
    \frac{\partial \mathcal{L}}{\partial \mathbf{A}} = s\mathbf{B}^\top \mathbf{G}^\top \mathbf{X}, \quad \frac{\partial \mathcal{L}}{\partial \mathbf{B}} = s\mathbf{G}^\top \mathbf{X}\mathbf{A}^\top,
    \label{eq:lora_grad}
\end{equation}
where $\mathbf{G} = \partial \mathcal{L} / \partial \mathbf{Y}$ is the gradient with respect to the output.

\begin{proposition}[Equivalent Weight Update for LoRA]
\label{prop:lora_update}
Given the LoRA parameterization in Eq.~\eqref{eq:lora_param}, after one SGD step with learning rate $\eta$, the effective weight change is:
\begin{equation}
    \Delta \mathbf{W}_{\text{LoRA}} = s(\mathbf{B}'\mathbf{A}' - \mathbf{B}\mathbf{A}) = -\eta s^2 \left( \mathbf{B}\mathbf{B}^\top \mathbf{G}^\top \mathbf{X} + \mathbf{G}^\top \mathbf{X}\mathbf{A}^\top \mathbf{A} \right) + O(\eta^2).
    \label{eq:lora_update}
\end{equation}
Under the assumption of small learning rate $\eta \ll 1$, the quadratic term $O(\eta^2)$ is negligible compared to the linear term $O(\eta)$.
\end{proposition}

The proof is provided in Appendix~\ref{app:prop1}.

\textbf{Null Space Projection as Degenerate LoRA.}
NSP mitigates catastrophic forgetting by constraining weight updates to directions that minimally affect historical representations. Consider a linear layer with output $\mathbf{Y} = \mathbf{X}\mathbf{W}^\top$, where $\mathbf{X} \in \mathbb{R}^{n \times d_{\text{in}}}$ is the input and $\mathbf{W} \in \mathbb{R}^{d_{\text{out}} \times d_{\text{in}}}$ the weight matrix. Under full-parameter fine-tuning, the weight gradient is $\mathbf{G}_{\mathbf{W}} = \mathbf{G}^\top \mathbf{X} \in \mathbb{R}^{d_{\text{out}} \times d_{\text{in}}}$, where $\mathbf{G} = \partial \mathcal{L} / \partial \mathbf{Y} \in \mathbb{R}^{n \times d_{\text{out}}}$. NSP projects this gradient onto the approximate null space of historical input features $\mathbf{X}_{\text{old}} \in \mathbb{R}^{n_{\text{old}} \times d_{\text{in}}}$, yielding the projected gradient $\mathbf{G}_{\mathbf{W}}^{\text{proj}} = \mathbf{G}_{\mathbf{W}} \mathbf{P}_{\text{null}}$ and the parameter update:
\begin{equation}
    \Delta \mathbf{W}_{\text{NSP}} = -\eta \, \mathbf{G}_{\mathbf{W}} \mathbf{P}_{\text{null}} = -\eta \, \mathbf{G}^\top \mathbf{X} \, \mathbf{P}_{\text{null}}.
    \label{eq:nsp_update}
\end{equation}
This ensures that $\mathbf{X}_{\text{old}} (\Delta \mathbf{W}_{\text{NSP}})^\top \approx \mathbf{0}$, i.e., the update does not alter outputs for historical inputs.

To approximate the null space incrementally, we maintain an online estimate of the feature Gram matrix. For task $t$, we collect input features $\mathbf{X}_t^{\text{raw}} \in \mathbb{R}^{n_t \times d_{\text{token}} \times d_{\text{in}}}$ from the LoRA layer, flatten them along the batch and token dimensions into $\mathbf{X}_t \in \mathbb{R}^{(n_t \cdot d_{\text{token}}) \times d_{\text{in}}}$, and compute the uncentered covariance (Gram matrix):
\begin{equation}
    \mathbf{\Sigma}_{\text{Cov}}^{(t)} = \alpha \, \mathbf{\Sigma}_{\text{Cov}}^{(t-1)} + (1-\alpha) \, \mathbf{X}_t^\top \mathbf{X}_t \in \mathbb{R}^{d_{\text{in}} \times d_{\text{in}}},
    \label{eq:cov_ema}
\end{equation}
where $\alpha \in [0,1)$ is a momentum coefficient. An eigendecomposition yields $\mathbf{\Sigma}_{\text{Cov}}^{(t)} = \mathbf{U}\mathbf{\Lambda} \mathbf{U}^\top$, where $\mathbf{U} = [\mathbf{u}_1, \ldots, \mathbf{u}_{d_{\text{in}}}]$ contains the eigenvectors and $\mathbf{\Lambda} = \mathrm{diag}(\lambda_1, \ldots, \lambda_{d_{\text{in}}})$ the eigenvalues sorted in descending order. The hard projection matrix onto the approximate null space is:
\begin{equation}
    \mathbf{P}_{\text{null}} = \mathbf{U}_2\mathbf{U}_2^\top = \mathbf{I}_{d_{\text{in}}} - \mathbf{U}_1\mathbf{U}_1^\top,
    \label{eq:hard_proj}
\end{equation}
where $\mathbf{U}_1 = [\mathbf{u}_1, \ldots, \mathbf{u}_k]$ spans the signal subspace and $\mathbf{U}_2 = [\mathbf{u}_{k+1}, \ldots, \mathbf{u}_{d_{\text{in}}}]$ the null subspace, with $k$ being the smallest integer satisfying $\sum_{i=1}^k \lambda_i \geq (1-\epsilon)\sum_{i=1}^{d_{\text{in}}} \lambda_i$ for a threshold $\epsilon > 0$. We use the uncentered Gram matrix because deep network features are typically zero-mean after batch normalization layers; the resulting eigenvectors approximately span the same subspace as the centered covariance.

\begin{theorem}[NSP as a Degenerate Form of LoRA]
\label{thm:nsp_lora}
Let $\mathbf{P}_{\text{null}} = \mathbf{V}_2\mathbf{V}_2^\top$ be the NSP projection matrix in Eq.~\eqref{eq:hard_proj}, where $\mathbf{V}_2 = \mathbf{U}_2$ spans the approximate null space of historical features. Consider a restricted LoRA parameterization where the input matrix is fixed as $\mathbf{A} = \frac{1}{s}\mathbf{V}_2^\top \in \mathbb{R}^{r \times d_{\text{in}}}$ (with $r = d_{\text{in}} - k$) and only the output matrix $\mathbf{B} \in \mathbb{R}^{d_{\text{out}} \times r}$ is trainable. Then the weight update produced by this restricted LoRA after one SGD step is identical to the NSP update in Eq.~\eqref{eq:nsp_update} (up to $O(\eta^2)$ terms).
\end{theorem}

The proof is provided in Appendix~\ref{app:thm1}.

\begin{definition}[Reachable Parameter Update Sets]
\label{def:reachable}
Consider a linear layer with pre-trained weight $\mathbf{W}_0 \in \mathbb{R}^{d_{\text{out}} \times d_{\text{in}}}$. For a parameterization of the form $\mathbf{W} = \mathbf{W}_0 + \Delta\mathbf{W}$, the \emph{reachable parameter update set} is the collection of all weight changes $\Delta\mathbf{W}$ achievable by varying the trainable parameters over their full domain:
\begin{itemize}
    \item \textbf{Standard LoRA:} $\mathcal{R}_{\text{LoRA}} = \{s\mathbf{B}\mathbf{A} : \mathbf{B} \in \mathbb{R}^{d_{\text{out}} \times r},\, \mathbf{A} \in \mathbb{R}^{r \times d_{\text{in}}}\}$;
    \item \textbf{NSP:} $\mathcal{R}_{\text{NSP}} = \{\mathbf{B}\mathbf{V}_2^\top : \mathbf{B} \in \mathbb{R}^{d_{\text{out}} \times (d_{\text{in}}-k)}\}$, where $\mathbf{V}_2$ spans the approximate null space as in Theorem~\ref{thm:nsp_lora};
    \item \textbf{LoRA-NSP:} $\mathcal{R}_{\text{LoRA-NSP}} = \{s\mathbf{B}\mathbf{A}\mathbf{P}_{\text{SGP}} : \mathbf{B} \in \mathbb{R}^{d_{\text{out}} \times r},\, \mathbf{A} \in \mathbb{R}^{r \times d_{\text{in}}}\}$.
\end{itemize}
\end{definition}

\begin{corollary}[Plasticity Limitation of NSP]
\label{cor:nsp_limit}
The reachable parameter update sets defined in Definition~\ref{def:reachable} satisfy $\mathcal{R}_{\text{NSP}} \subset \mathcal{R}_{\text{LoRA}}$, i.e., NSP's reachable set is a strict subset of LoRA's.
\end{corollary}

\begin{proof}[Proof Sketch]
From Theorem~\ref{thm:nsp_lora}, NSP is equivalent to a restricted LoRA where $\mathbf{A}$ is fixed to $\frac{1}{s}\mathbf{V}_2^\top$. In standard LoRA, both $\mathbf{A}$ and $\mathbf{B}$ are free to vary over $\mathbb{R}^{r \times d_{\text{in}}}$ and $\mathbb{R}^{d_{\text{out}} \times r}$, respectively. The restricted parameterization constrains $\mathbf{A}$ to a fixed subspace, eliminating all signal-subspace directions from the update. Thus any update achievable by NSP is also achievable by LoRA (by setting $\mathbf{A}$ to the same fixed value), but the converse does not hold since standard LoRA can vary $\mathbf{A}$ over the full space.
\end{proof}

\textbf{Soft Projected LoRA.}
Theorem~\ref{thm:nsp_lora} and Corollary~\ref{cor:nsp_limit} reveal that NSP is a degenerate form of LoRA with severely limited plasticity: by fixing $\mathbf{A}$ to span only the null space, the effective rank of the update is reduced from $r$ to $d_{\text{in}} - k$, and all signal-subspace directions are completely blocked. This explains NSP's poor adaptability in practice. Motivated by this insight, we propose a \textit{soft forward-pass projection} mechanism that retains all eigendirections with adaptive weighting, improving the stability-plasticity trade-off while preserving LoRA's full expressivity.

The key insight is that the forward-pass projection $\mathbf{P}_{\text{SGP}}$ modulates the effective weight update by reweighting the contribution of each input eigendirection. From Eq.~\eqref{eq:lora_nsp_grad}, both the $\mathbf{A}$ and $\mathbf{B}$ gradients are scaled by $\mathbf{P}_{\text{SGP}}$ from the right, which modulates how each input direction contributes to the parameter update. Under the standard PEFT assumption that LoRA updates are small relative to the frozen backbone ($s\mathbf{B}\mathbf{A}\mathbf{P}_{\text{SGP}} \ll \mathbf{W}_0$), the coupling between layers induced by backpropagation is of higher order and can be neglected in the first-order analysis. Consequently, the forward-pass parameterization induces a modulation effect on the gradient directions that is approximately equivalent to applying a soft projection directly to the LoRA gradients.

Given the covariance matrix $\mathbf{\Sigma}_{\text{Cov}}^{(t)}$ from Eq.~\eqref{eq:cov_ema}, let $\mathbf{\Sigma}_{\text{Cov}}^{(t)} = \mathbf{U}\mathbf{\Lambda} \mathbf{U}^\top$ be its eigendecomposition. We define the soft projection matrix as:
\begin{equation}
    \mathbf{P}_{\text{SGP}}(\mathbf{\Lambda}, \mathbf{U}) = \sum_{i=1}^{d_{\text{in}}} w(\lambda_i) \, \mathbf{u}_i \mathbf{u}_i^\top,
    \label{eq:psgp}
\end{equation}
where $w(\lambda) \in (0,1]$ is a monotonically decreasing weighting function. We employ:
\begin{equation}
    w(\lambda) = \frac{1}{1 + \beta\ln(1 + \lambda^p)},
    \label{eq:weight_fn}
\end{equation}
where $\beta > 0$ controls the decay rate and $p > 0$ the sensitivity to eigenvalue magnitude. Since $\mathbf{P}_{\text{SGP}}$ is a linear combination of symmetric rank-1 matrices with positive coefficients, it is symmetric positive definite when $w(\lambda_i) > 0$ for all $i$.

In LoRA-NSP, we introduce this soft projection into the low-rank update. The adapted weight becomes:
\begin{equation}
    \mathbf{W} = \mathbf{W}_0 + s\mathbf{B}\mathbf{A}\mathbf{P}_{\text{SGP}}.
    \label{eq:lora_nsp_fwd}
\end{equation}
Since $\mathbf{P}_{\text{SGP}}$ acts in the forward pass, it naturally modulates the gradient flow during backpropagation. The gradients with respect to the trainable matrices are:
\begin{equation}
    \frac{\partial \mathcal{L}}{\partial \mathbf{A}} = s\mathbf{B}^\top \mathbf{G}^\top \mathbf{X} \mathbf{P}_{\text{SGP}}, \quad \frac{\partial \mathcal{L}}{\partial \mathbf{B}} = s\mathbf{G}^\top \mathbf{X} \mathbf{P}_{\text{SGP}} \mathbf{A}^\top,
    \label{eq:lora_nsp_grad}
\end{equation}
which shows that the projection matrix scales both the $\mathbf{A}$ and $\mathbf{B}$ gradient directions consistently. This forward-pass formulation offers better compatibility with existing PEFT techniques and optimizers than backward-pass gradient modification methods like GPM.

\begin{theorem}[Expressive Equivalence of LoRA-NSP and LoRA]
\label{thm:expressive}
Let $\mathcal{R}_{\text{LoRA}}$ and $\mathcal{R}_{\text{LoRA-NSP}}$ be the reachable parameter update sets defined in Definition~\ref{def:reachable}. If all weighting coefficients satisfy $w(\lambda_i) > 0$ and $\mathbf{\Sigma}_{\text{Cov}}^{(t)}$ is full-rank, then $\mathbf{P}_{\text{SGP}}$ is invertible and $\mathcal{R}_{\text{LoRA-NSP}} = \mathcal{R}_{\text{LoRA}}$.
\end{theorem}

The proof is provided in Appendix~\ref{app:thm2}.

Theorem~\ref{thm:expressive} establishes that LoRA-NSP unifies LoRA and NSP: by choosing $w(\lambda)$, LoRA-NSP smoothly interpolates between standard LoRA ($w(\lambda) \approx 1$, so $\mathbf{P}_{\text{SGP}} \approx \mathbf{I}$) and hard NSP ($w(\lambda) \to 0$ for signal directions, so $\mathbf{P}_{\text{SGP}} \approx \mathbf{P}_{\text{null}}$).

\textbf{Discussion on Multi-Layer Coupling.}
The gradient expressions in Eq.~\eqref{eq:lora_nsp_grad} hold for each LoRA layer individually. In a deep network where multiple layers employ LoRA-NSP simultaneously, a subtlety arises: the backward-propagated gradient $\mathbf{G} = \partial \mathcal{L} / \partial \mathbf{Y}$ for layer $l$ depends on the parameters of all subsequent layers $l+1, \ldots, L$. Consequently, the effective gradient for layer $l$ is modulated not only by its own $\mathbf{P}_{\text{SGP}}^{(l)}$ but also indirectly by $\mathbf{P}_{\text{SGP}}^{(l+1)}, \ldots, \mathbf{P}_{\text{SGP}}^{(L)}$ through the backpropagation chain.

However, under the standard PEFT assumption that LoRA updates are small ($s\mathbf{B}\mathbf{A}\mathbf{P}_{\text{SGP}} \ll \mathbf{W}_0$), this cross-layer coupling is a higher-order effect. To see this, consider a two-layer network with LoRA-NSP in both layers. The backward gradient for layer 1 is $\mathbf{G}_1 = \mathbf{G}_2 \mathbf{W}_2 = \mathbf{G}_2 (\mathbf{W}_{2,0} + s\mathbf{B}_2\mathbf{A}_2\mathbf{P}_{\text{SGP}}^{(2)})$. The term $\mathbf{G}_2 \mathbf{W}_{2,0}$ is the ``frozen-backbone'' gradient, while $\mathbf{G}_2 (s\mathbf{B}_2\mathbf{A}_2\mathbf{P}_{\text{SGP}}^{(2)})$ is the coupling correction. Since $s\mathbf{B}_2\mathbf{A}_2\mathbf{P}_{\text{SGP}}^{(2)} = O(s)$, the correction is $O(s)$ relative to the backbone gradient. The gradient for layer 1's LoRA parameters thus receives a first-order modulation from $\mathbf{P}_{\text{SGP}}^{(1)}$ and a second-order correction from $\mathbf{P}_{\text{SGP}}^{(2)}$. In the small-update regime ($s \to 0$), the first-order term dominates, and the per-layer analysis in Eq.~\eqref{eq:lora_nsp_grad} remains accurate. Our empirical validation on the X-TAIL benchmark confirms that this approximation is well-justified in practice.

\textbf{Covariance History and Incremental Updates.}
For each LoRA layer, we maintain the covariance history using the exponential moving average in Eq.~\eqref{eq:cov_ema}. After eigendecomposition $\mathbf{\Sigma}_{\text{Cov}}^{(t)} = \mathbf{U} \mathbf{\Lambda} \mathbf{U}^\top$, we construct the soft projection matrix via Eq.~\eqref{eq:psgp}--\eqref{eq:weight_fn}.

\textbf{Incremental Learning with Weight Merging.}
After training on each task, we merge the LoRA weights into the frozen backbone via $\mathbf{W}_0 \leftarrow \mathbf{W}_0 + s\mathbf{B}\mathbf{A}\mathbf{P}_{\text{SGP}}$, then reinitialize $\mathbf{A} \sim \mathcal{N}(\mathbf{0},\mathbf{I})$ and reset $\mathbf{B} \leftarrow \mathbf{0}$, preparing for the next task. This prevents parameter growth while accumulating knowledge across tasks.

\textbf{FD+CD Distillation Loss.}
To further maintain representation stability, we design a distillation loss combining feature distillation (FD) and cross-modal distillation (CD) on a reference dataset (Flickr8K):
\begin{enumerate}
    \item \textbf{Feature Distillation (FD):} Aligns image features between the fine-tuned model and the pre-trained model:
    \begin{equation}
        \mathcal{L}_{\text{fd}} = \|f_I^{\text{pre}}(\mathbf{x}) - f_I^{\text{student}}(\mathbf{x})\|_2^2
    \end{equation}
    \item \textbf{Cross-modal Distillation (CD):} Preserves vision-language alignment using contrastive loss with temperature $\tau$:
    \begin{equation}
        \mathcal{L}_{\text{cd}} = -\frac{1}{N}\sum_{i=1}^N \log \frac{\exp(s_{ii}/\tau)}{\sum_{j=1}^N \exp(s_{ij}/\tau)}
    \end{equation}
    where $s_{ij}$ denotes the cosine similarity between image $i$ and text $j$ features.
\end{enumerate}
The total training loss is:
\begin{equation}
    \mathcal{L} = \mathcal{L}_{\text{CE}} + \lambda_{\text{fd}} \mathcal{L}_{\text{fd}} + \lambda_{\text{cd}} \mathcal{L}_{\text{cd}}
\end{equation}

\subsection{Inference Stage: Ensemble Classifier}

\textbf{Motivation.}
After fine-tuning, the zero-shot classifier leverages CLIP's pre-trained vision-language alignment for classification. However, it relies solely on visual-text similarity and may miss discriminative information present in the adapted visual feature space. We study an ensemble classifier that combines a Low-Rank Regularized Gaussian Discriminant Analysis (LR-RGDA) classifier with the zero-shot classifier. The classifier is motivated by Gaussian discriminant analysis and is implemented with low-rank covariance corrections for computational efficiency.

\textbf{Regularized Gaussian Discriminant Analysis.}
Building upon Gaussian statistics, we construct a classifier that approximates the Gaussian discriminant decision rule. After training on task $t$, we model the deep features of each class $c \in \mathcal{C}_t$ as a multivariate Gaussian distribution $p(\mathbf{z}|c) = \mathcal{N}(\boldsymbol{\mu}_c, \mathbf{\Sigma}_c)$, where $\mathbf{z} = f_\theta(\mathbf{x})$ is the feature representation and $\mathcal{C}_t$ is the set of classes observed up to task $t$.

By Bayes' theorem, the posterior satisfies $p(c|\mathbf{x}) \propto p(c) \, p(\mathbf{z}|c)$, where $p(c)$ is the class prior. Under uniform priors $p(c) = 1/|\mathcal{C}_t|$, the log-posterior (dropping class-independent constants) is:
\begin{equation}
    \log p(c|\mathbf{x}) \propto -\frac{1}{2}(\mathbf{z} - \boldsymbol{\mu}_c)^\top \mathbf{\Sigma}_c^{-1} (\mathbf{z} - \boldsymbol{\mu}_c) - \frac{1}{2}\log|\mathbf{\Sigma}_c| + \log p(c).
    \label{eq:log_posterior}
\end{equation}

In practice, the class-wise covariance $\mathbf{\Sigma}_c$ estimated from limited samples is often ill-conditioned or singular. We propose a composite regularization scheme:
\begin{equation}
    \mathbf{\Sigma}_c^{\text{reg}} = \alpha_1 \mathbf{\Sigma}_c + \alpha_2 \mathbf{\Sigma}_{\text{avg}} + \alpha_3 \mathbf{I}_d,
    \label{eq:cov_reg}
\end{equation}
where $\mathbf{\Sigma}_{\text{avg}} = \frac{1}{|\mathcal{C}_t|}\sum_{c \in \mathcal{C}_t} \mathbf{\Sigma}_c$ captures shared feature correlations across classes, $\mathbf{I}_d$ ensures numerical stability, and $\alpha_1, \alpha_2, \alpha_3 \geq 0$ control the contribution of each component. Setting $\alpha_1 = 1, \alpha_2 = \alpha_3 = 0$ recovers standard Quadratic Discriminant Analysis (QDA), whereas $\alpha_1 = \alpha_3 = 0, \alpha_2 = 1$ reduces to Linear Discriminant Analysis (LDA).

Substituting Eq.~\eqref{eq:cov_reg} into Eq.~\eqref{eq:log_posterior} yields the RGDA discriminant function:
\begin{equation}
    g_c^{\text{RGDA}}(\mathbf{z}) = -\frac{1}{2}(\mathbf{z} - \boldsymbol{\mu}_c)^\top (\mathbf{\Sigma}_c^{\text{reg}})^{-1} (\mathbf{z} - \boldsymbol{\mu}_c) - \frac{1}{2}\log|\mathbf{\Sigma}_c^{\text{reg}}| + \log p(c).
    \label{eq:rgda}
\end{equation}
The RGDA classifier selects the class with maximum discriminant score: $\hat{y} = \arg\max_{c \in \mathcal{C}_t} g_c^{\text{RGDA}}(\mathbf{z})$. Under correctly specified Gaussian class-conditionals and the chosen regularized covariance model, this rule corresponds to the plug-in Gaussian discriminant classifier discussed in Appendix~\ref{app:bayes}.

\textbf{Low-Rank Factorized RGDA.}
The full RGDA inference complexity scales as $\mathcal{O}(Cd^2)$ due to per-class matrix-vector products with $(\mathbf{\Sigma}_c^{\text{reg}})^{-1}$. This becomes expensive when both the number of classes $C$ and the feature dimension $d$ are large.

To address this, we exploit the observation that class covariance matrices typically exhibit low-rank structures around a shared baseline. We decompose the regularized covariance as:
\begin{equation}
    \mathbf{\Sigma}_c^{\text{reg}} = \mathbf{B} + \mathbf{U}_c \boldsymbol{\Lambda}_c \mathbf{U}_c^\top,
    \label{eq:lr_cov}
\end{equation}
where $\mathbf{B} = \alpha_2 \mathbf{\Sigma}_{\text{avg}} + \alpha_3 \mathbf{I}_d$ is a global full-rank base matrix shared across all classes, $\mathbf{U}_c \in \mathbb{R}^{d \times r}$ contains the top-$r$ eigenvectors of $(\mathbf{\Sigma}_c^{\text{reg}} - \mathbf{B})$, and $\boldsymbol{\Lambda}_c = \mathrm{diag}(\lambda_{c,1}, \ldots, \lambda_{c,r}) \in \mathbb{R}^{r \times r}$ stores the corresponding eigenvalues, with $r \ll d$. Notably, the low-rank perturbation captures class-specific deviations from the global covariance structure, enabling fine-grained class discrimination without incurring full-rank costs.

Applying the Woodbury matrix identity to Eq.~\eqref{eq:lr_cov} yields:
\begin{equation}
    (\mathbf{\Sigma}_c^{\text{reg}})^{-1} = \mathbf{B}^{-1} - \mathbf{B}^{-1} \mathbf{U}_c \big(\boldsymbol{\Lambda}_c^{-1} + \mathbf{U}_c^\top \mathbf{B}^{-1} \mathbf{U}_c\big)^{-1} \mathbf{U}_c^\top \mathbf{B}^{-1}.
    \label{eq:woodbury}
\end{equation}
Since $\mathbf{B}$ is shared across classes, its inverse is computed once and cached. The term $(\boldsymbol{\Lambda}_c^{-1} + \mathbf{U}_c^\top \mathbf{B}^{-1} \mathbf{U}_c)^{-1}$ involves only $r \times r$ matrices, reducing the per-sample per-class cost from $\mathcal{O}(d^2)$ to $\mathcal{O}(dr + r^3)$. With $r \ll d$, the total inference complexity becomes $\mathcal{O}(Cd + Cr^2)$, which is comparable to linear classifiers (see Appendix~\ref{app:complexity} for a detailed derivation).

Substituting Eq.~\eqref{eq:woodbury} into Eq.~\eqref{eq:rgda} and expanding the quadratic form, the LR-RGDA discriminant function decomposes into a global linear term plus a class-specific quadratic correction in a rank-$r$ subspace:
\begin{equation}
    g_c^{\text{LR-RGDA}}(\mathbf{z}) = \mathbf{w}_c^\top \mathbf{z} + b_c + \mathbf{z}^\top \mathbf{Q}_c \mathbf{z},
    \label{eq:lr_rgda}
\end{equation}
where
\begin{align}
    \mathbf{w}_c &= \mathbf{B}^{-1} \boldsymbol{\mu}_c - \mathbf{B}^{-1} \mathbf{U}_c \big(\boldsymbol{\Lambda}_c^{-1} + \mathbf{U}_c^\top \mathbf{B}^{-1} \mathbf{U}_c\big)^{-1} \mathbf{U}_c^\top \mathbf{B}^{-1} \boldsymbol{\mu}_c, \label{eq:wc}\\
    b_c &= -\frac{1}{2} \boldsymbol{\mu}_c^\top \mathbf{B}^{-1} \boldsymbol{\mu}_c + \frac{1}{2} \boldsymbol{\mu}_c^\top \mathbf{B}^{-1} \mathbf{U}_c \big(\boldsymbol{\Lambda}_c^{-1} + \mathbf{U}_c^\top \mathbf{B}^{-1} \mathbf{U}_c\big)^{-1} \mathbf{U}_c^\top \mathbf{B}^{-1} \boldsymbol{\mu}_c \nonumber\\
    &\quad -\frac{1}{2}\log|\mathbf{B} + \mathbf{U}_c \boldsymbol{\Lambda}_c \mathbf{U}_c^\top| + \log p(c), \label{eq:bc}\\
    \mathbf{Q}_c &= -\frac{1}{2} \mathbf{B}^{-1} \mathbf{U}_c \big(\boldsymbol{\Lambda}_c^{-1} + \mathbf{U}_c^\top \mathbf{B}^{-1} \mathbf{U}_c\big)^{-1} \mathbf{U}_c^\top \mathbf{B}^{-1}. \label{eq:qc}
\end{align}
The complete derivation from Eq.~\eqref{eq:woodbury} to Eq.~\eqref{eq:lr_rgda} is provided in Appendix~\ref{app:lr_rgda_derivation}.

\textbf{Ensemble Framework.}
The ensemble classifier balances between supervised classification from LR-RGDA and CLIP's zero-shot classifier through a contribution coefficient $\alpha$:
\begin{equation}
    p_{\text{ens}}(y|\mathbf{x}) = \alpha \cdot p_{\text{LR-RGDA}}(y|\mathbf{x}) + (1-\alpha) \cdot p_{\text{ZS}}(y|\mathbf{x})
\end{equation}
over the evaluation label space. In our current implementation we use a fixed mixture rather than an OOD router. This avoids introducing a separate threshold-selection problem and keeps the comparison with LADA focused on classifier construction from the same feature source. We report real-feature and statistical-replay variants separately because they answer different memory-budget questions.

%==============================================================================
\fi
\section{Experiments}
\label{sec:experiments}

\subsection{Experimental Setup}

\textbf{Datasets.}
We conduct experiments on the recently proposed X-TAIL~\cite{xu2024rail} benchmark, which consists of 10 image classification datasets: Aircraft~\cite{maji2013fine}, Caltech101~\cite{fei2004learning}, DTD~\cite{cimpoi2014describing}, EuroSAT~\cite{helber2019eurosat}, Flowers~\cite{nilsback2008automated}, Food~\cite{bossard2014food}, MNIST~\cite{deng2012mnist}, OxfordPet~\cite{parkhi2012cats}, StanfordCars~\cite{krause2013collecting}, and SUN397~\cite{xiao2010sun}. Each dataset is treated as a task, and the benchmark includes a total of 1,100 classes across the 10 tasks. Following prior work~\cite{zheng2023zscl}, we evaluate under the 16-shot setting.

\textbf{Evaluation Metrics.}
Following prior work~\cite{zheng2023zscl}, we employ three metrics:
\begin{itemize}
    \item \textbf{Transfer:} Average accuracy on tasks $k+1, k+2, ..., K$ after training on task $k$, quantifying forward forgetting.
    \item \textbf{Average:} Mean accuracy across all time steps, offering a holistic measure of stability and plasticity.
    \item \textbf{Last:} Model performance after continual training on all tasks, capturing plasticity and backward forgetting.
\end{itemize}

\textbf{Baselines.}
We compare against the following methods:
\begin{itemize}
    \item \textbf{Zero-shot:} Vanilla CLIP without fine-tuning.
    \item \textbf{LwF~\cite{li2017learning}:} Learning without Forgetting using knowledge distillation.
    \item \textbf{WiSE-FT~\cite{wortsman2022robust}:} Robust fine-tuning via weight ensemble.
    \item \textbf{ZSCL~\cite{zheng2023zscl}:} Zero-shot transfer degradation prevention via regularization.
    \item \textbf{MoE-Adapters~\cite{yu2024moe}:} Mixture-of-experts adapters for continual learning.
    \item \textbf{RAIL~\cite{xu2024rail}:} Regression-based analytic incremental learning.
    \item \textbf{LADA~\cite{luo2025lada}:} Label-specific CLIP adapter.
\end{itemize}

\textbf{Implementation Details.}
We adopt the CLIP~\cite{radford2021learning} model with a ViT-B/16~\cite{dosovitskiy2021an} image encoder. The training process uses AdamW~\cite{loshchilov2019decoupled} optimizer with learning rate $10^{-4}$ and weight decay $3 \times 10^{-5}$. For LoRA-NSP, we set rank $r=4$, NSP weight $\lambda_{\text{nsp}}=0.02$, and covariance momentum $\alpha=0.9$. For the ensemble classifier, we set $\alpha=0.5$ and temperature $T=1.0$. For LR-RGDA, we use rank 32 with regularization coefficients $\alpha_1=0.6$, $\alpha_2=1.0$, $\alpha_3=0.5$.

\subsection{Main Results}

Strict task-incremental Transfer/Average/Last results are not yet used as paper evidence in this draft. We keep the LADA and RAIL reference numbers as protocol targets, but final claims about LoRA-NSP will be added only after the planned three-seed training-side ablation and incremental evaluation are complete. The current verified result is the storage-dependent classifier boundary in Section~\ref{sec:stat_replay}.

\iffalse

\begin{table}[t]
\centering
\caption{Comparison of different continual learning methods on X-TAIL 16-shot setting for each task in terms of Transfer, Average, and Last scores (\%). The best results are highlighted with \textbf{bold} style.}
\label{tab:main_16shot}
\resizebox{\textwidth}{!}{
\begin{tabular}{l|cccccccccc|c}
\toprule
Method & \rotatebox{90}{Aircraft} & \rotatebox{90}{Caltech101} & \rotatebox{90}{DTD} & \rotatebox{90}{EuroSAT} & \rotatebox{90}{Flowers} & \rotatebox{90}{Food} & \rotatebox{90}{MNIST} & \rotatebox{90}{Pets} & \rotatebox{90}{Cars} & \rotatebox{90}{SUN397} & Average \\
\midrule
Zero-shot & 23.8 & 74.3 & 36.4 & 37.4 & 64.1 & 83.4 & 43.9 & 87.8 & 65.5 & 60.8 & 57.7 \\
\midrule
\multicolumn{12}{l}{\textbf{Transfer}} \\
LwF~\cite{li2017learning} & --- & --- & --- & --- & --- & --- & --- & --- & --- & --- & 47.7 \\
WiSE-FT~\cite{wortsman2022robust} & --- & --- & --- & --- & --- & --- & --- & --- & --- & --- & 52.3 \\
ZSCL~\cite{zheng2023zscl} & --- & --- & --- & --- & --- & --- & --- & --- & --- & --- & 59.0 \\
MoE-Adapters~\cite{yu2024moe} & --- & --- & --- & --- & --- & --- & --- & --- & --- & --- & 56.0 \\
Primal-RAIL~\cite{xu2024rail} & --- & --- & --- & --- & --- & --- & --- & --- & --- & --- & 71.1 \\
Dual-RAIL~\cite{xu2024rail} & --- & --- & --- & --- & --- & --- & --- & --- & --- & --- & 71.3 \\
LADA~\cite{luo2025lada} & --- & --- & --- & --- & --- & --- & --- & --- & --- & --- & 61.5 \\
\textbf{Ours (LoRA-NSP)} & \todo{--} & \todo{--} & \todo{--} & \todo{--} & \todo{--} & \todo{--} & \todo{--} & \todo{--} & \todo{--} & \todo{--} & \todo{--} \\
\midrule
\multicolumn{12}{l}{\textbf{Average}} \\
LwF~\cite{li2017learning} & --- & --- & --- & --- & --- & --- & --- & --- & --- & --- & 53.2 \\
WiSE-FT~\cite{wortsman2022robust} & --- & --- & --- & --- & --- & --- & --- & --- & --- & --- & 54.2 \\
ZSCL~\cite{zheng2023zscl} & --- & --- & --- & --- & --- & --- & --- & --- & --- & --- & 60.0 \\
MoE-Adapters~\cite{yu2024moe} & --- & --- & --- & --- & --- & --- & --- & --- & --- & --- & 63.0 \\
Primal-RAIL~\cite{xu2024rail} & --- & --- & --- & --- & --- & --- & --- & --- & --- & --- & 71.3 \\
Dual-RAIL~\cite{xu2024rail} & --- & --- & --- & --- & --- & --- & --- & --- & --- & --- & 71.3 \\
LADA~\cite{luo2025lada} & --- & --- & --- & --- & --- & --- & --- & --- & --- & --- & 72.7 \\
\textbf{Ours (LoRA-NSP)} & \todo{--} & \todo{--} & \todo{--} & \todo{--} & \todo{--} & \todo{--} & \todo{--} & \todo{--} & \todo{--} & \todo{--} & \todo{--} \\
\midrule
\multicolumn{12}{l}{\textbf{Last}} \\
LwF~\cite{li2017learning} & --- & --- & --- & --- & --- & --- & --- & --- & --- & --- & 62.8 \\
WiSE-FT~\cite{wortsman2022robust} & --- & --- & --- & --- & --- & --- & --- & --- & --- & --- & 58.0 \\
ZSCL~\cite{zheng2023zscl} & --- & --- & --- & --- & --- & --- & --- & --- & --- & --- & 63.4 \\
MoE-Adapters~\cite{yu2024moe} & --- & --- & --- & --- & --- & --- & --- & --- & --- & --- & 70.5 \\
Primal-RAIL~\cite{xu2024rail} & --- & --- & --- & --- & --- & --- & --- & --- & --- & --- & 81.4 \\
Dual-RAIL~\cite{xu2024rail} & --- & --- & --- & --- & --- & --- & --- & --- & --- & --- & 82.3 \\
LADA~\cite{luo2025lada} & --- & --- & --- & --- & --- & --- & --- & --- & --- & --- & 83.1 \\
\textbf{Ours (LoRA-NSP)} & \todo{--} & \todo{--} & \todo{--} & \todo{--} & \todo{--} & \todo{--} & \todo{--} & \todo{--} & \todo{--} & \todo{--} & \todo{--} \\
\bottomrule
\end{tabular}
}
\end{table}

\textbf{16-shot Setting.}
\todo{Fill narrative based on actual results. Our method (LoRA-NSP with FD+CD distillation) achieves [Transfer]\% Transfer, [Average]\% Average, and [Last]\% Last accuracy. Compare with zero-shot baseline and other methods in Table~\ref{tab:main_16shot}. Discuss trade-offs with RAIL and other baselines.}

\todo{[Key Observations] Remove or revise based on actual results:
\begin{itemize}
    \item Our training-stage optimization (LoRA-NSP + FD+CD distillation) provides consistent improvements over the zero-shot baseline.
    \item Methods that completely freeze CLIP features (RAIL) achieve high Transfer but lower Average/Last scores, while our approach balances all three metrics.
    \item The null space parameterization helps maintain representation stability across tasks.
\end{itemize}
}

\subsection{Ablation Study}

\begin{table}[t]
\centering
\caption{Ablation study on 5-task X-TAIL setting. We evaluate contributions of LoRA baseline, null space parameterization (NSP), feature distillation (FD), and cross-modal distillation (CD). All values are percentages.}
\label{tab:ablation}
\begin{tabular}{lccc|ccc}
\toprule
LoRA & NSP & FD & CD & Transfer & Average & Last \\
\midrule
\checkmark & & & & \todo{50.9} & \todo{54.1} & \todo{64.4} \\
\checkmark & \checkmark & & & \todo{56.1} & \todo{60.4} & \todo{72.4} \\
\checkmark & & \checkmark & & \todo{56.1} & \todo{60.4} & \todo{72.4} \\
\checkmark & & & \checkmark & \todo{56.3} & \todo{61.7} & \todo{74.4} \\
\checkmark & \checkmark & \checkmark & \checkmark & \todo{62.5} & \todo{68.3} & \todo{74.5} \\
\bottomrule
\end{tabular}
\end{table}

\todo{[Ablation Narrative] Fill based on actual ablation results. Expected structure:
\begin{itemize}
    \item \textbf{NSP:} Adding null space projection improves Transfer from [X] to [Y] ([±Z] pp).
    \item \textbf{FD:} Feature distillation contributes [±X] pp on Transfer and [±Y] pp on Average.
    \item \textbf{CD:} Cross-modal distillation contributes [±X] pp on Transfer and [±Y] pp on Average.
    \item \textbf{Full combination:} Achieves [Transfer]\% Transfer, [Avg]\% Average, and [Last]\% Last.
\end{itemize}
}

\subsection{Analysis of Ensemble Classifier}

The ensemble classifier combines the LR-RGDA classifier (trained on seen classes) with the zero-shot classifier through a contribution coefficient $\alpha$. We evaluate the classifier variants on the full 10-task X-TAIL benchmark. Results are reported in Table~\ref{tab:alpha}.

\begin{table}[t]
\centering
\caption{Comparison of classifier variants on X-TAIL 16-shot. We evaluate zero-shot classifier, LR-RGDA classifier, and their ensemble ($\alpha=0.5$).}
\label{tab:alpha}
\begin{tabular}{l|ccc}
\toprule
Classifier & Transfer & Average & Last \\
\midrule
Zero-shot & --- & 57.7 & --- \\
LR-RGDA only & \todo{48.5} & \todo{65.2} & \todo{78.3} \\
\textbf{Ensemble ($\alpha=0.5$)} & \todo{63.0} & \todo{65.6} & \todo{70.2} \\
\bottomrule
\end{tabular}
\end{table}

\todo{[Classifier Analysis Narrative] Fill based on actual results. Expected content:
\begin{itemize}
    \item Zero-shot classifier: [characteristic description].
    \item LR-RGDA only: [characteristic description with actual numbers].
    \item Ensemble: [characteristic description with actual numbers].
\end{itemize}
}
\fi

\subsection{Statistical Replay Protocol}
\label{sec:stat_replay}

To isolate the inference-stage contribution, we evaluate classifier construction under two feature-storage protocols. In the \emph{real-feature} protocol, classifiers are fitted from the deterministic 16-shot feature set for each class. In the \emph{statistical-replay} protocol, previous class features are represented by compact Gaussian statistics and replayed through class means or GMM-derived pseudo-features. Both LR-RGDA and LADA are fitted from the same feature source, use the same deterministic classifier transform, and are evaluated on the same full test set.

This protocol is intentionally narrower than LADA's official DPT training procedure. LADA DPT uses historical prototypes or Gaussian samples as replay during incremental training, while the current experiment freezes the feature source and rebuilds the final classifiers from matched real or statistical features. Therefore, the comparison below should be read as a controlled classifier-construction study under equal replay sources, not as a claim that compact GMM classifier rebuilding is equivalent to official LADA DPT. This distinction lets us test the classifier inductive bias under a fixed memory budget without conflating it with training-time replay dynamics.

\begin{table}[t]
\centering
\caption{Classifier comparison under real-feature and statistical-replay protocols on X-TAIL 16-shot. All classifiers are built with \texttt{classifier\_feature\_transform=test} and evaluated on the same full test set. Results are averaged over seeds 42/43/44.}
\label{tab:gmm_mean_replay}
\resizebox{\textwidth}{!}{%
\begin{tabular}{lccccc}
\toprule
Replay protocol & CLIP-ZS & LR-RGDA & LADA & LR-RGDA+ZS & LADA+ZS \\
\midrule
Real 16-shot features & $56.40{\pm}0.00$ & $77.87{\pm}0.16$ & $\mathbf{78.92{\pm}0.16}$ & $78.05{\pm}0.14$ & $\mathbf{78.92{\pm}0.16}$ \\
GMM raw component means & $56.40{\pm}0.00$ & $76.93{\pm}0.16$ & $76.08{\pm}0.35$ & $\mathbf{77.12{\pm}0.13}$ & $76.08{\pm}0.35$ \\
GMM raw samples & $56.40{\pm}0.00$ & $76.24{\pm}0.11$ & $75.63{\pm}0.28$ & $\mathbf{76.46{\pm}0.16}$ & $75.63{\pm}0.28$ \\
GMM sphere samples & $56.40{\pm}0.00$ & $76.22{\pm}0.18$ & $75.66{\pm}0.31$ & $\mathbf{76.45{\pm}0.17}$ & $75.66{\pm}0.31$ \\
\bottomrule
\end{tabular}
}
\end{table}

Table~\ref{tab:gmm_mean_replay} shows a clear storage-dependent boundary. When real 16-shot features are available, LADA remains stronger ($78.92{\pm}0.16$) than LR-RGDA ($77.87{\pm}0.16$) and LR-RGDA+ZS ($78.05{\pm}0.14$). Under compact statistical replay, however, LR-RGDA consistently exceeds LADA. The strongest setting is GMM component-mean replay: across the three seeds, the LR-RGDA - LADA margins are $+0.96$, $+0.99$, and $+0.59$ percentage points, while the LR-RGDA+ZS - LADA margins are $+1.18$, $+1.19$, and $+0.74$ points. Stochastic GMM replay gives smaller but still positive margins for both raw-space and sphere-space sampling, suggesting that the main effect is not the raw-versus-sphere fitting space but the use of compact, low-noise component statistics. We therefore interpret LR-RGDA as a statistical-replay classifier that is most competitive when historical features must be compressed, not as a universal replacement for LADA under unrestricted feature storage.

\begin{table}[t]
\centering
\caption{Replay-source storage budget for X-TAIL 16-shot with CLIP ViT-B/16 features ($C=1100$ classes, $d=512$, float32 storage). The table counts the stored historical feature representation used to construct or replay classifiers, not the final classifier parameters.}
\label{tab:storage_budget}
\begin{tabular}{lccc}
\toprule
Stored representation & Items per class & Formula & Storage \\
\midrule
Real 16-shot features & 16 features & $C \cdot 16 \cdot d$ & 34.38 MiB \\
LADA $k{=}16$ centers & 16 centers & $C \cdot 16 \cdot d$ & 34.38 MiB \\
GMM $k{=}4$ component means & 4 means & $C \cdot 4 \cdot d$ & 8.59 MiB \\
GMM $k{=}4$ spherical params & 4 means, variances, weights & $C \cdot 4 \cdot (d+2)$ & 8.63 MiB \\
\bottomrule
\end{tabular}
\end{table}

Table~\ref{tab:storage_budget} makes the memory setting explicit. With 16-shot data, storing LADA's $k=16$ label centers has the same feature-vector count as storing all real training features. By contrast, the $k=4$ GMM replay source used in our compact statistical protocol stores roughly one quarter as many feature vectors. This is the memory regime in which LR-RGDA is most competitive in Table~\ref{tab:gmm_mean_replay}. We do not claim that the final LR-RGDA classifier parameters are always smaller than LADA's label-memory classifier; the claim is about reconstructing or fitting classifiers from a compact historical replay source.

% OOD detection experiments were part of an earlier draft. They are intentionally
% removed from the main claim until a dedicated, leakage-free evaluation is added.
\iffalse
\subsection{OOD Detection Performance}

We evaluate the OOD detection capability of our LR-RGDA-based detector. Table~\ref{tab:ood} reports the AUROC and FPR@95TPR metrics on X-TAIL.

\begin{table}[t]
\centering
\caption{OOD detection performance on X-TAIL. We report AUROC and FPR@95 (lower is better).}
\label{tab:ood}
\begin{tabular}{lcc}
\toprule
Method & AUROC ($\uparrow$) & FPR@95 ($\downarrow$) \\
\midrule
Mahalanobis~\cite{lee2018simple} & 82.3 & 45.2 \\
Energy~\cite{liu2020energy} & 84.1 & 42.8 \\
MSP~\cite{hendrycks2017baseline} & 79.5 & 52.1 \\
\midrule
\textbf{Ours (LR-RGDA)} & \todo{AUROC} & \todo{FPR@95} \\
\bottomrule
\end{tabular}
\end{table}

\todo{[OOD Detection Narrative] Fill based on actual results. Our LR-RGDA-based OOD detector achieves [AUROC]\% AUROC and [FPR]\% FPR@95, [comparison with baselines].}
\fi

%==============================================================================
\section{Conclusion}
\label{sec:conclusion}

In this paper, we propose a two-stage framework for CLIP continual learning that combines LoRA-NSP training with inference-time classifier construction. At the training stage, LoRA-NSP integrates soft null-space constraints with low-rank adaptation to mitigate catastrophic forgetting while retaining plasticity. At the inference stage, the LR-RGDA and zero-shot ensemble provides a statistical classifier alternative to label-memory baselines, with the strongest current motivation in compact statistical replay settings. Extensive experiments on the X-TAIL benchmark are ongoing; final claims will be restricted to settings where three-seed results verify an advantage over LADA or demonstrate a clear memory-performance trade-off.

%==============================================================================
\bibliographystyle{plain}
\begin{thebibliography}{10}

\bibitem{radford2021learning}
A.~Radford et al.
\newblock Learning transferable visual models from natural language supervision.
\newblock In {\em ICML}, 2021.

\bibitem{wang2024lora}
Y.~Wang et al.
\newblock LoRA-CL: Continual learning with low-rank adaptation.
\newblock In {\em ICLR}, 2024.

\bibitem{french1999catastrophic}
R.~M. French.
\newblock Catastrophic forgetting in connectionist networks.
\newblock {\em Trends in Cognitive Sciences}, 1999.

\bibitem{mermillod2013stability}
M.~Mermillod et al.
\newblock The stability-plasticity dilemma: Investigating the continuum from catastrophic forgetting to aged-limited learning effects.
\newblock {\em Frontiers in Psychology}, 2013.

\bibitem{xu2024rail}
X.~Xu et al.
\newblock RAIL: Regression-based analytic incremental learning for cross-domain task-agnostic learning.
\newblock In {\em NeurIPS}, 2024.

\bibitem{hsu2018re}
Y.-C. Hsu et al.
\newblock Re-evaluating continual learning scenarios: A categorization and case for strong baselines.
\newblock In {\em NeurIPS Continual Learning Workshop}, 2018.

\bibitem{rebuffi2017icarl}
S.-A. Rebuffi et al.
\newblock iCaRL: Incremental classifier and representation learning.
\newblock In {\em CVPR}, 2017.

\bibitem{zheng2023zscl}
Y.~Zheng et al.
\newblock ZSCL: Preventing zero-shot transfer degradation in continual learning of vision-language models.
\newblock In {\em CVPR}, 2023.

\bibitem{de2021continual}
M.~De~Lange et al.
\newblock A continual learning survey: Defying forgetting in classification tasks.
\newblock {\em TPAMI}, 2021.

\bibitem{wang2024b}
Y.~Wang et al.
\newblock A comprehensive survey of continual learning: Theory, method and application.
\newblock {\em TPAMI}, 2024.

\bibitem{kirkpatrick2017overcoming}
J.~Kirkpatrick et al.
\newblock Overcoming catastrophic forgetting in neural networks.
\newblock {\em PNAS}, 2017.

\bibitem{ahn2019uncertainty}
H.~Ahn et al.
\newblock Uncertainty-based continual learning with adaptive regularization.
\newblock In {\em NeurIPS}, 2019.

\bibitem{rolnick2019experience}
D.~Rolnick et al.
\newblock Experience replay for continual learning.
\newblock In {\em NeurIPS}, 2019.

\bibitem{wang2022learning}
Z.~Wang et al.
\newblock Learning to prompt for continual learning.
\newblock In {\em CVPR}, 2022.

\bibitem{wang2022dualprompt}
Z.~Wang et al.
\newblock DualPrompt: Complementary prompting for rehearsal-free continual learning.
\newblock In {\em ECCV}, 2022.

\bibitem{wang2022sprompts}
Z.~Wang et al.
\newblock S-Prompts Learning with pre-trained transformers: An Occam's razor for domain incremental learning.
\newblock In {\em NeurIPS}, 2022.

\bibitem{yu2024moe}
J.~Yu et al.
\newblock MoE-Adapters: Mixture-of-experts adapters for parameter-efficient continual learning.
\newblock In {\em CVPR}, 2024.

\bibitem{luo2025lada}
M.-L. Luo et al.
\newblock LADA: Scalable label-specific CLIP adapter for continual learning.
\newblock In {\em ICML}, 2025.

\bibitem{ni2023modx}
Z. Ni et al.
\newblock Continual vision-language representation learning with off-diagonal information.
\newblock In {\em ICML}, 2023.

\bibitem{huang2025mgclip}
L. Huang et al.
\newblock Mind the gap: Preserving and compensating for the modality gap in CLIP-based continual learning.
\newblock In {\em ICCV}, 2025.

\bibitem{liu2025cclip}
W. Liu et al.
\newblock C-CLIP: Multimodal continual learning for vision-language model.
\newblock In {\em ICLR}, 2025.

\bibitem{hu2022lora}
E.~J. Hu et al.
\newblock LoRA: Low-rank adaptation of large language models.
\newblock In {\em ICLR}, 2022.

\bibitem{li2017learning}
Z.~Li and D.~Hoiem.
\newblock Learning without forgetting.
\newblock {\em TPAMI}, 2017.

\bibitem{wortsman2022robust}
M.~Wortsman et al.
\newblock Robust fine-tuning of zero-shot models.
\newblock In {\em NeurIPS}, 2022.

\bibitem{dosovitskiy2021an}
A.~Dosovitskiy et al.
\newblock An image is worth 16x16 words: Transformers for image recognition at scale.
\newblock In {\em ICLR}, 2021.

\bibitem{loshchilov2019decoupled}
I.~Loshchilov and F.~Hutter.
\newblock Decoupled weight decay regularization.
\newblock In {\em ICLR}, 2019.

\bibitem{maji2013fine}
S.~Maji et al.
\newblock Fine-grained visual classification of aircraft.
\newblock {\em arXiv:1306.5151}, 2013.

\bibitem{fei2004learning}
L.~Fei-Fei et al.
\newblock Learning generative visual models from few training examples.
\newblock In {\em CVPR Workshops}, 2004.

\bibitem{cimpoi2014describing}
M.~Cimpoi et al.
\newblock Describing textures in the wild.
\newblock In {\em CVPR}, 2014.

\bibitem{helber2019eurosat}
P.~Helber et al.
\newblock EuroSAT: A novel dataset and deep learning benchmark for land use and land cover classification.
\newblock {\em IEEE JSTARS}, 2019.

\bibitem{nilsback2008automated}
M.-E. Nilsback and A.~Zisserman.
\newblock Automated flower classification over a large number of classes.
\newblock In {\em ICVGIP}, 2008.

\bibitem{bossard2014food}
L.~Bossard et al.
\newblock Food-101 -- mining discriminative components with random forests.
\newblock In {\em ECCV}, 2014.

\bibitem{deng2012mnist}
L.~Deng.
\newblock The MNIST database of handwritten digit images for machine learning research.
\newblock {\em IEEE Signal Processing Magazine}, 2012.

\bibitem{parkhi2012cats}
O.~M. Parkhi et al.
\newblock Cats and dogs.
\newblock In {\em CVPR}, 2012.

\bibitem{krause2013collecting}
J.~Krause et al.
\newblock Collecting a large-scale dataset of fine-grained cars.
\newblock In {\em FGVC Workshop}, 2013.

\bibitem{xiao2010sun}
J.~Xiao et al.
\newblock SUN database: Large-scale scene recognition from abbey to zoo.
\newblock In {\em CVPR}, 2010.

\bibitem{lee2018simple}
K.~Lee et al.
\newblock A simple unified framework for detecting out-of-distribution samples and adversarial attacks.
\newblock In {\em NeurIPS}, 2018.

\bibitem{liu2020energy}
W.~Liu et al.
\newblock Energy-based out-of-distribution detection.
\newblock In {\em NeurIPS}, 2020.

\bibitem{hendrycks2017baseline}
D.~Hendrycks and K.~Gimpel.
\newblock A baseline for detecting misclassified and out-of-distribution examples in neural networks.
\newblock In {\em ICLR}, 2017.

\bibitem{saha2021gradient}
G.~Saha et al.
\newblock Gradient projection memory for continual learning.
\newblock In {\em ICLR}, 2021.

\bibitem{kang2025dynamic}
B.~Kang et al.
\newblock Dynamic multi-layer null space projection for vision-language continual learning.
\newblock In {\em ICCV}, 2025.

\bibitem{tang2026lora}
P.~Tang et al.
\newblock Put the space of LoRA initialization to the extreme to preserve pre-trained knowledge.
\newblock In {\em AAAI}, 2026.

\bibitem{lu2026codyra}
H.~Lu et al.
\newblock Adaptive rank, reduced forgetting: Knowledge retention in continual learning vision-language models with dynamic rank-selective LoRA.
\newblock In {\em ICLR}, 2026.

\end{thebibliography}

%==============================================================================
\appendix
%==============================================================================

\section{MAP Equivalence of RGDA Under a Regularized Gaussian Model}
\label{app:bayes}

\begin{theorem}[MAP Equivalence of RGDA]
\label{thm:bayes}
Let $\mathcal{C}_t = \{1, \ldots, C\}$ be the set of classes observed up to task $t$. Assume that:
\begin{enumerate}[label=(\roman*)]
    \item The class-conditional feature distribution is Gaussian: $p(\mathbf{z}|c) = \mathcal{N}(\boldsymbol{\mu}_c, \mathbf{\Sigma}_c)$ for each $c \in \mathcal{C}_t$;
    \item Class priors are uniform: $p(c) = 1/C$ for all $c \in \mathcal{C}_t$;
    \item The regularized covariance matrices $\mathbf{\Sigma}_c^{\text{reg}}$ defined in Eq.~\eqref{eq:cov_reg} are positive definite.
\end{enumerate}
Then the RGDA classifier $\hat{y} = \arg\max_{c \in \mathcal{C}_t} g_c^{\text{RGDA}}(\mathbf{z})$ is equivalent to the maximum a posteriori (MAP) decision rule under the regularized Gaussian model $p^{\text{reg}}(\mathbf{z}|c) = \mathcal{N}(\boldsymbol{\mu}_c, \mathbf{\Sigma}_c^{\text{reg}})$.
\end{theorem}

\begin{proof}
Under the regularized Gaussian model $p^{\text{reg}}(\mathbf{z}|c) = \mathcal{N}(\boldsymbol{\mu}_c, \mathbf{\Sigma}_c^{\text{reg}})$, the class-conditional log-likelihood is:
\begin{equation}
    \log p^{\text{reg}}(\mathbf{z}|c) = -\frac{1}{2}(\mathbf{z} - \boldsymbol{\mu}_c)^\top (\mathbf{\Sigma}_c^{\text{reg}})^{-1} (\mathbf{z} - \boldsymbol{\mu}_c) - \frac{1}{2}\log|\mathbf{\Sigma}_c^{\text{reg}}| - \frac{d}{2}\log(2\pi).
\end{equation}
By Bayes' theorem, the posterior log-probability satisfies:
\begin{equation}
    \log p^{\text{reg}}(c|\mathbf{z}) = \log p^{\text{reg}}(\mathbf{z}|c) + \log p(c) - \log p(\mathbf{z}).
\end{equation}
Since $\log p(\mathbf{z})$ is class-independent, maximizing the posterior is equivalent to maximizing $\log p^{\text{reg}}(\mathbf{z}|c) + \log p(c)$. Substituting the Gaussian log-likelihood and dropping the constant $-\frac{d}{2}\log(2\pi)$ yields:
\begin{equation}
    \log p^{\text{reg}}(c|\mathbf{z}) \;\propto\; -\frac{1}{2}(\mathbf{z} - \boldsymbol{\mu}_c)^\top (\mathbf{\Sigma}_c^{\text{reg}})^{-1} (\mathbf{z} - \boldsymbol{\mu}_c) - \frac{1}{2}\log|\mathbf{\Sigma}_c^{\text{reg}}| + \log p(c) = g_c^{\text{RGDA}}(\mathbf{z}).
\end{equation}
Thus $\arg\max_{c} g_c^{\text{RGDA}}(\mathbf{z}) = \arg\max_{c} p^{\text{reg}}(c|\mathbf{z})$, completing the proof.
\end{proof}

%==============================================================================
\section{LR-RGDA Discriminant Function Derivation}
\label{app:lr_rgda_derivation}

\begin{proposition}[Decomposition of LR-RGDA Discriminant Function]
\label{prop:decomposition}
Let $\mathbf{\Sigma}_c^{\text{reg}} = \mathbf{B} + \mathbf{U}_c \boldsymbol{\Lambda}_c \mathbf{U}_c^\top$ as defined in Eq.~\eqref{eq:lr_cov}, with $\mathbf{B} \succ 0$. Then the RGDA discriminant function $g_c^{\text{RGDA}}(\mathbf{z})$ in Eq.~\eqref{eq:rgda} can be decomposed as $g_c^{\text{LR-RGDA}}(\mathbf{z}) = \mathbf{w}_c^\top \mathbf{z} + b_c + \mathbf{z}^\top \mathbf{Q}_c \mathbf{z}$, where $\mathbf{w}_c$, $b_c$, and $\mathbf{Q}_c$ are given by Eqs.~\eqref{eq:wc}--\eqref{eq:qc}.
\end{proposition}

\begin{proof}
We begin from the Woodbury identity (Eq.~\eqref{eq:woodbury}):
\begin{equation}
    (\mathbf{\Sigma}_c^{\text{reg}})^{-1} = \mathbf{B}^{-1} - \mathbf{B}^{-1} \mathbf{U}_c \mathbf{M}_c^{-1} \mathbf{U}_c^\top \mathbf{B}^{-1},
\end{equation}
where we abbreviate $\mathbf{M}_c = \boldsymbol{\Lambda}_c^{-1} + \mathbf{U}_c^\top \mathbf{B}^{-1} \mathbf{U}_c$ for notational brevity. Let $\boldsymbol{\delta}_c = \mathbf{z} - \boldsymbol{\mu}_c$. Substituting into the quadratic form of Eq.~\eqref{eq:rgda}:
\begin{align}
    \boldsymbol{\delta}_c^\top (\mathbf{\Sigma}_c^{\text{reg}})^{-1} \boldsymbol{\delta}_c
    &= \boldsymbol{\delta}_c^\top \mathbf{B}^{-1} \boldsymbol{\delta}_c - \boldsymbol{\delta}_c^\top \mathbf{B}^{-1} \mathbf{U}_c \mathbf{M}_c^{-1} \mathbf{U}_c^\top \mathbf{B}^{-1} \boldsymbol{\delta}_c \nonumber\\
    &= (\mathbf{z}^\top \mathbf{B}^{-1} \mathbf{z} - 2\boldsymbol{\mu}_c^\top \mathbf{B}^{-1} \mathbf{z} + \boldsymbol{\mu}_c^\top \mathbf{B}^{-1} \boldsymbol{\mu}_c) \nonumber\\
    &\quad - \bigl(
    \mathbf{z}^\top \mathbf{B}^{-1} \mathbf{U}_c \mathbf{M}_c^{-1}
    \mathbf{U}_c^\top \mathbf{B}^{-1} \mathbf{z} \nonumber\\
    &\qquad
    - 2\boldsymbol{\mu}_c^\top \mathbf{B}^{-1} \mathbf{U}_c \mathbf{M}_c^{-1}
    \mathbf{U}_c^\top \mathbf{B}^{-1} \mathbf{z} \nonumber\\
    &\qquad
    + \boldsymbol{\mu}_c^\top \mathbf{B}^{-1} \mathbf{U}_c \mathbf{M}_c^{-1}
    \mathbf{U}_c^\top \mathbf{B}^{-1} \boldsymbol{\mu}_c
    \bigr). \label{eq:quad_expand}
\end{align}
Collecting terms by powers of $\mathbf{z}$:
\begin{itemize}
    \item \textbf{Quadratic term in $\mathbf{z}$:} $\mathbf{z}^\top \mathbf{B}^{-1} \mathbf{z} - \mathbf{z}^\top \mathbf{B}^{-1} \mathbf{U}_c \mathbf{M}_c^{-1} \mathbf{U}_c^\top \mathbf{B}^{-1} \mathbf{z} = \mathbf{z}^\top (\mathbf{\Sigma}_c^{\text{reg}})^{-1} \mathbf{z}$.
    Since $g_c^{\text{RGDA}}$ carries a factor of $-\frac{1}{2}$, this contributes $\mathbf{z}^\top \mathbf{Q}_c \mathbf{z}$ with:
    \begin{equation}
        \mathbf{Q}_c = -\frac{1}{2}(\mathbf{\Sigma}_c^{\text{reg}})^{-1} = -\frac{1}{2}\mathbf{B}^{-1} \mathbf{U}_c \mathbf{M}_c^{-1} \mathbf{U}_c^\top \mathbf{B}^{-1},
    \end{equation}
    which matches Eq.~\eqref{eq:qc} after absorbing the global $\mathbf{B}^{-1}$ term into the linear part.

    \item \textbf{Linear term in $\mathbf{z}$:} From Eq.~\eqref{eq:quad_expand}, the coefficients of $\mathbf{z}$ (excluding the quadratic self-term) are:
    \begin{equation}
        2\boldsymbol{\mu}_c^\top \mathbf{B}^{-1} \mathbf{z} - 2\boldsymbol{\mu}_c^\top \mathbf{B}^{-1} \mathbf{U}_c \mathbf{M}_c^{-1} \mathbf{U}_c^\top \mathbf{B}^{-1} \mathbf{z} = 2\big(\mathbf{B}^{-1} \boldsymbol{\mu}_c - \mathbf{B}^{-1} \mathbf{U}_c \mathbf{M}_c^{-1} \mathbf{U}_c^\top \mathbf{B}^{-1} \boldsymbol{\mu}_c\big)^\top \mathbf{z}.
    \end{equation}
    Multiplying by $-\frac{1}{2}$ gives the linear term $\mathbf{w}_c^\top \mathbf{z}$ with $\mathbf{w}_c$ as in Eq.~\eqref{eq:wc}.

    \item \textbf{Constant term $b_c$:} Collecting all terms independent of $\mathbf{z}$ from Eq.~\eqref{eq:quad_expand} and multiplying by $-\frac{1}{2}$:
    \begin{align}
        b_c &= -\frac{1}{2}\boldsymbol{\mu}_c^\top \mathbf{B}^{-1} \boldsymbol{\mu}_c + \frac{1}{2}\boldsymbol{\mu}_c^\top \mathbf{B}^{-1} \mathbf{U}_c \mathbf{M}_c^{-1} \mathbf{U}_c^\top \mathbf{B}^{-1} \boldsymbol{\mu}_c \nonumber\\
        &\quad -\frac{1}{2}\log|\mathbf{B} + \mathbf{U}_c \boldsymbol{\Lambda}_c \mathbf{U}_c^\top| + \log p(c),
    \end{align}
    which is exactly Eq.~\eqref{eq:bc}.
\end{itemize}
\end{proof}

%==============================================================================
\section{Complexity Analysis of LR-RGDA}
\label{app:complexity}

We provide a detailed complexity analysis for LR-RGDA inference, comparing it against RGDA and a standard linear classifier. Let $C$ denote the number of classes, $d$ the feature dimension, and $r$ the low-rank dimension ($r \ll d$). All matrices are assumed real-valued.

\subsection{Inference Complexity per Sample}

\textbf{RGDA.} For each class $c$, computing $g_c^{\text{RGDA}}(\mathbf{z})$ requires:
\begin{enumerate}
    \item One matrix-vector product: $(\mathbf{\Sigma}_c^{\text{reg}})^{-1} (\mathbf{z} - \boldsymbol{\mu}_c)$, costing $\mathcal{O}(d^2)$.
    \item One inner product with $(\mathbf{z} - \boldsymbol{\mu}_c)$, costing $\mathcal{O}(d)$.
    \item The log-determinant term is precomputed offline.
\end{enumerate}
The total per-sample complexity is $\mathcal{O}(Cd^2)$.

\textbf{LR-RGDA.} Using the decomposition in Eq.~\eqref{eq:lr_rgda}, the per-class computation proceeds as follows:
\begin{enumerate}
    \item Compute $\mathbf{v}_c = \mathbf{U}_c^\top \mathbf{B}^{-1} \mathbf{z}$. Since $\mathbf{B}^{-1}$ is precomputed and cached, this costs $\mathcal{O}(dr)$ (matrix-vector with $\mathbf{B}^{-1}$) plus $\mathcal{O}(dr)$ (matrix-vector with $\mathbf{U}_c^\top$), totaling $\mathcal{O}(dr)$.
    \item Solve the $r \times r$ linear system $\mathbf{M}_c \mathbf{u}_c = \mathbf{v}_c$ where $\mathbf{M}_c = \boldsymbol{\Lambda}_c^{-1} + \mathbf{U}_c^\top \mathbf{B}^{-1} \mathbf{U}_c$. This costs $\mathcal{O}(r^3)$ via Cholesky decomposition (since $\mathbf{M}_c \succ 0$) or $\mathcal{O}(r^2)$ if $\mathbf{M}_c^{-1}$ is precomputed.
    \item Compute the quadratic correction $\mathbf{z}^\top \mathbf{Q}_c \mathbf{z} = -\frac{1}{2} \mathbf{v}_c^\top \mathbf{u}_c$ via an inner product costing $\mathcal{O}(r)$.
    \item Compute the linear term $\mathbf{w}_c^\top \mathbf{z}$. Since $\mathbf{w}_c$ is precomputed (Eq.~\eqref{eq:wc}), this costs $\mathcal{O}(d)$.
\end{enumerate}
Summing across all $C$ classes and assuming precomputation of $\mathbf{M}_c^{-1}$, the total per-sample complexity is:
\begin{equation}
    \mathcal{O}\big(C(dr + r^2 + d)\big) = \mathcal{O}(Cd + Cr^2),
\end{equation}
where the $Cr^2$ term dominates when $r^2 > d$, and the $Cd$ term dominates otherwise.

\textbf{Linear Classifier.} A standard linear classifier computes $C$ dot products $\mathbf{w}_c^\top \mathbf{z} + b_c$, costing $\mathcal{O}(Cd)$ per sample.

\subsection{Precomputation Costs}

\begin{table}[h]
\centering
\caption{Precomputation and per-sample inference complexity comparison.}
\label{tab:complexity}
\begin{tabular}{lcc}
\toprule
\textbf{Component} & \textbf{Precomputation} & \textbf{Per-Sample Inference} \\
\midrule
RGDA & $\mathcal{O}(Cd^3)$ (inverse + det) & $\mathcal{O}(Cd^2)$ \\
LR-RGDA & $\mathcal{O}(Cd^3 + Cr^3)$ (eigen-decomp + $\mathbf{M}_c^{-1}$) & $\mathcal{O}(Cd + Cr^2)$ \\
Linear Classifier & $\mathcal{O}(Cd)$ (weight storage) & $\mathcal{O}(Cd)$ \\
\bottomrule
\end{tabular}
\end{table}

Table~\ref{tab:complexity} summarizes the complexity. The precomputation for LR-RGDA includes:
\begin{itemize}
    \item Computing $\mathbf{B}^{-1}$: $\mathcal{O}(d^3)$ once.
    \item Eigen-decomposition of $(\mathbf{\Sigma}_c^{\text{reg}} - \mathbf{B})$ for each class: $\mathcal{O}(Cd^3)$.
    \item Computing and inverting $\mathbf{M}_c$ for each class: $\mathcal{O}(Cr^3)$.
    \item Computing $\mathbf{w}_c$ and $b_c$ for each class: $\mathcal{O}(Cdr)$.
\end{itemize}
Since precomputation is performed once after training on each task (not per sample), it does not affect the online inference latency.

\subsection{Storage Complexity}

For $C$ classes, the storage requirements are:
\begin{itemize}
    \item \textbf{RGDA:} $C$ mean vectors $\boldsymbol{\mu}_c \in \mathbb{R}^d$ and $C$ full covariance matrices (or their inverses) $\in \mathbb{R}^{d \times d}$, totaling $\mathcal{O}(Cd + Cd^2) = \mathcal{O}(Cd^2)$.
    \item \textbf{LR-RGDA:} $C$ mean vectors, $C$ matrices $\mathbf{U}_c \in \mathbb{R}^{d \times r}$, $C$ diagonal matrices $\boldsymbol{\Lambda}_c \in \mathbb{R}^{r \times r}$, and the shared $\mathbf{B}^{-1} \in \mathbb{R}^{d \times d}$, totaling $\mathcal{O}(Cd + Cdr + Cr + d^2) = \mathcal{O}(d^2 + Cdr)$.
    \item \textbf{Linear Classifier:} $C$ weight vectors and biases, totaling $\mathcal{O}(Cd)$.
\end{itemize}
When $r \ll d$, LR-RGDA reduces storage from $\mathcal{O}(Cd^2)$ to $\mathcal{O}(d^2 + Cdr)$, a substantial saving for large $C$ and $d$.

\subsection{Practical Considerations}

In our experiments with CLIP ViT-B/16 ($d = 512$) and $r = 16$, LR-RGDA achieves:
\begin{itemize}
    \item \textbf{Inference speedup:} $Cd^2 / (Cd + Cr^2) \approx d/r^2 = 512/256 = 2\times$ faster than full RGDA when $C$ is large.
    \item \textbf{Storage reduction:} $Cd^2 / (d^2 + Cdr) \approx d/r = 32\times$ for $C \gg d/r$.
\end{itemize}
Moreover, the per-sample operations (matrix-vector products with $\mathbf{B}^{-1}$ and $\mathbf{U}_c$) are highly parallelizable on modern hardware, enabling further speedups through GPU batching.

%==============================================================================
\section{Proofs for LoRA-NF}
\label{app:lora_nf_proofs}

\subsection{Proof of Proposition~\ref{prop:historical_response}}

Using the definition of the Frobenius norm and cyclic invariance of the trace,
\begin{align}
    \frac{1}{N_h}\|\mathbf{X}_h\mathbf{P}\|_F^2
    &=\frac{1}{N_h}\operatorname{tr}
      (\mathbf{P}^\top\mathbf{X}_h^\top\mathbf{X}_h\mathbf{P})\\
    &=\operatorname{tr}(\mathbf{P}^\top\mathbf{G}_h\mathbf{P}).
\end{align}
When $\mathbf{P}=\mathbf{U}\operatorname{diag}(p_i)\mathbf{U}^\top$ and
$\mathbf{G}_h=\mathbf{U}\operatorname{diag}(\lambda_i)\mathbf{U}^\top$, the trace reduces to $\sum_i p_i^2\lambda_i$. The second statement follows from submultiplicativity:
\begin{equation}
    \|\mathbf{X}_h\mathbf{P}\mathbf{A}\mathbf{B}\|_F
    \leq\|\mathbf{X}_h\mathbf{P}\|_F\|\mathbf{A}\|_2\|\mathbf{B}\|_2.
\end{equation}

\subsection{Proof of Proposition~\ref{prop:expressive_equivalence}}

For any LoRA-NF update, define $\mathbf{C}=\mathbf{P}\mathbf{A}$; then
$\mathbf{P}\mathbf{A}\mathbf{B}=\mathbf{C}\mathbf{B}$, proving inclusion in the rank-matched LoRA set. Conversely, for any $\mathbf{C}\mathbf{B}$, invertibility of $\mathbf{P}$ permits $\mathbf{A}=\mathbf{P}^{-1}\mathbf{C}$, for which $\mathbf{P}\mathbf{A}\mathbf{B}=\mathbf{C}\mathbf{B}$. The two sets are therefore equal. For Eq.~\eqref{eq:leaky_filter}, the eigenvalues of $\mathbf{P}$ are $1$ and $\rho$; hence $\rho>0$ is sufficient for invertibility.

\subsection{Proof of Proposition~\ref{prop:optimization_geometry}}

Because $\mathbf{C}=\mathbf{P}\mathbf{A}$ and $\mathbf{P}$ is fixed,
\begin{equation}
    \nabla_{\mathbf{A}}\mathcal{L}
    =\mathbf{P}^\top\nabla_{\mathbf{C}}\mathcal{L}.
\end{equation}
A gradient-descent step on $\mathbf{A}$ gives
\begin{align}
    \mathbf{C}^{+}
    &=\mathbf{P}\left(\mathbf{A}-\eta\nabla_{\mathbf{A}}\mathcal{L}\right)\\
    &=\mathbf{C}-\eta\mathbf{P}\mathbf{P}^\top
      \nabla_{\mathbf{C}}\mathcal{L}.
\end{align}
If $\mathbf{P}$ is symmetric, $\mathbf{P}\mathbf{P}^\top=\mathbf{P}^2$. This result concerns the induced geometry of gradient descent; adaptive optimizers additionally transform the parameter-space gradients through their moment estimates.

\iffalse
\section{Proof of Theorem~\ref{thm:nsp_lora}: NSP as a Degenerate Form of LoRA}
\label{app:thm1}

\begin{proof}[Proof of Theorem~\ref{thm:nsp_lora}]
Consider the restricted LoRA parameterization where $\mathbf{A} = \frac{1}{s}\mathbf{V}_2^\top \in \mathbb{R}^{r \times d_{\text{in}}}$ is fixed with $r = d_{\text{in}} - k$, and only $\mathbf{B} \in \mathbb{R}^{d_{\text{out}} \times r}$ is trainable. Under this parameterization, the adapted weight matrix is:
\begin{equation}
    \mathbf{W} = \mathbf{W}_0 + s\mathbf{B}\mathbf{A} = \mathbf{W}_0 + s\mathbf{B} \left(\frac{1}{s}\mathbf{V}_2^\top\right) = \mathbf{W}_0 + \mathbf{B}\mathbf{V}_2^\top.
\end{equation}

Since $\mathbf{A}$ is fixed, only $\mathbf{B}$ receives gradients. From Eq.~\eqref{eq:lora_grad}, the gradient with respect to $\mathbf{B}$ is:
\begin{equation}
    \frac{\partial \mathcal{L}}{\partial \mathbf{B}} = s\mathbf{G}^\top \mathbf{X} \mathbf{A}^\top = s\mathbf{G}^\top \mathbf{X} \left(\frac{1}{s}\mathbf{V}_2\right) = \mathbf{G}^\top \mathbf{X} \mathbf{V}_2.
\end{equation}

After one SGD step with learning rate $\eta$, the updated matrix is:
\begin{equation}
    \mathbf{B}' = \mathbf{B} - \eta \, \mathbf{G}^\top \mathbf{X} \mathbf{V}_2.
\end{equation}

The effective weight change is:
\begin{align}
    \Delta \mathbf{W}_{\text{restricted}}
    &= s(\mathbf{B}'\mathbf{A} - \mathbf{B}\mathbf{A}) \\
    &= s\left[\left(\mathbf{B} - \eta \, \mathbf{G}^\top \mathbf{X} \mathbf{V}_2\right)\left(\frac{1}{s}\mathbf{V}_2^\top\right) - \mathbf{B}\left(\frac{1}{s}\mathbf{V}_2^\top\right)\right] \\
    &= \left(\mathbf{B} - \eta \, \mathbf{G}^\top \mathbf{X} \mathbf{V}_2\right)\mathbf{V}_2^\top - \mathbf{B}\mathbf{V}_2^\top \\
    &= -\eta \, \mathbf{G}^\top \mathbf{X} \mathbf{V}_2 \mathbf{V}_2^\top \\
    &= -\eta \, \mathbf{G}^\top \mathbf{X} \, \mathbf{P}_{\text{null}}.
\end{align}

This is exactly the NSP update in Eq.~\eqref{eq:nsp_update}. Note that unlike standard LoRA where both $\mathbf{A}$ and $\mathbf{B}$ are trainable, here $\mathbf{A}$ is fixed, so there is no $O(\eta^2)$ term arising from the product of $\mathbf{A}$ and $\mathbf{B}$ updates; the equivalence is exact after a single SGD step.
\end{proof}

%==============================================================================
\section{Proof of Theorem~\ref{thm:expressive}: Expressive Equivalence of LoRA-NSP and LoRA}
\label{app:thm2}

\begin{proof}[Proof of Theorem~\ref{thm:expressive}]
Recall that $\mathbf{P}_{\text{SGP}} = \sum_{i=1}^{d_{\text{in}}} w(\lambda_i) \, \mathbf{u}_i \mathbf{u}_i^\top = \mathbf{U} \, \mathrm{diag}\big(w(\lambda_1), \ldots, w(\lambda_{d_{\text{in}}})\big) \, \mathbf{U}^\top$, where $\mathbf{U}$ is the orthogonal eigenvector matrix of $\mathbf{\Sigma}_{\text{Cov}}^{(t)}$.

\textbf{Invertibility of $\mathbf{P}_{\text{SGP}}$.} Since $w(\lambda_i) > 0$ for all $i$ by assumption, the diagonal matrix $\mathbf{D} = \mathrm{diag}\big(w(\lambda_1), \ldots, w(\lambda_{d_{\text{in}}})\big)$ is positive definite. Because $\mathbf{U}$ is orthogonal, we have:
\begin{equation}
    \mathbf{P}_{\text{SGP}} = \mathbf{U}\mathbf{D}\mathbf{U}^\top \succ 0,
\end{equation}
which implies $\mathbf{P}_{\text{SGP}}$ is symmetric positive definite and therefore invertible, with $\mathbf{P}_{\text{SGP}}^{-1} = \mathbf{U}\mathbf{D}^{-1}\mathbf{U}^\top$.

\textbf{Set equality.} Let $\mathcal{R}_{\text{LoRA-NSP}} = \{s\mathbf{B}\mathbf{A}\mathbf{P}_{\text{SGP}} : \mathbf{B} \in \mathbb{R}^{d_{\text{out}} \times r}, \mathbf{A} \in \mathbb{R}^{r \times d_{\text{in}}}\}$ and $\mathcal{R}_{\text{LoRA}} = \{s\mathbf{B}\mathbf{A} : \mathbf{B} \in \mathbb{R}^{d_{\text{out}} \times r}, \mathbf{A} \in \mathbb{R}^{r \times d_{\text{in}}}\}$.

$(\subseteq)$ For any $\mathbf{M} \in \mathcal{R}_{\text{LoRA-NSP}}$, there exist $\mathbf{B}, \mathbf{A}$ such that $\mathbf{M} = s\mathbf{B}\mathbf{A}\mathbf{P}_{\text{SGP}}$. Define $\mathbf{A}' = \mathbf{A}\mathbf{P}_{\text{SGP}} \in \mathbb{R}^{r \times d_{\text{in}}}$. Then $\mathbf{M} = s\mathbf{B}\mathbf{A}' \in \mathcal{R}_{\text{LoRA}}$. Hence $\mathcal{R}_{\text{LoRA-NSP}} \subseteq \mathcal{R}_{\text{LoRA}}$.

$(\supseteq)$ For any $\mathbf{M} \in \mathcal{R}_{\text{LoRA}}$, there exist $\mathbf{B}, \mathbf{A}$ such that $\mathbf{M} = s\mathbf{B}\mathbf{A}$. Define $\mathbf{A}' = \mathbf{A}\mathbf{P}_{\text{SGP}}^{-1} \in \mathbb{R}^{r \times d_{\text{in}}}$, which is well-defined because $\mathbf{P}_{\text{SGP}}$ is invertible. Then:
\begin{equation}
    s\mathbf{B}\mathbf{A}'\mathbf{P}_{\text{SGP}} = s\mathbf{B}\big(\mathbf{A}\mathbf{P}_{\text{SGP}}^{-1}\big)\mathbf{P}_{\text{SGP}} = s\mathbf{B}\mathbf{A} = \mathbf{M}.
\end{equation}
Thus $\mathbf{M} \in \mathcal{R}_{\text{LoRA-NSP}}$, proving $\mathcal{R}_{\text{LoRA}} \subseteq \mathcal{R}_{\text{LoRA-NSP}}$.

Combining both directions yields $\mathcal{R}_{\text{LoRA-NSP}} = \mathcal{R}_{\text{LoRA}}$.
\end{proof}

%==============================================================================
\section{Proof of Proposition~\ref{prop:lora_update}: Equivalent Weight Update for LoRA}
\label{app:prop1}

\begin{proof}[Proof of Proposition~\ref{prop:lora_update}]
Under the LoRA parameterization $\mathbf{W} = \mathbf{W}_0 + s\mathbf{B}\mathbf{A}$, the trainable parameters are $\mathbf{A} \in \mathbb{R}^{r \times d_{\text{in}}}$ and $\mathbf{B} \in \mathbb{R}^{d_{\text{out}} \times r}$. After one SGD step with learning rate $\eta$, the updated matrices are:
\begin{equation}
    \mathbf{A}' = \mathbf{A} - \eta \frac{\partial \mathcal{L}}{\partial \mathbf{A}}, \quad \mathbf{B}' = \mathbf{B} - \eta \frac{\partial \mathcal{L}}{\partial \mathbf{B}}.
\end{equation}
Substituting the gradients from Eq.~\eqref{eq:lora_grad}:
\begin{equation}
    \mathbf{A}' = \mathbf{A} - \eta s \mathbf{B}^\top \mathbf{G}^\top \mathbf{X}, \quad \mathbf{B}' = \mathbf{B} - \eta s \mathbf{G}^\top \mathbf{X}\mathbf{A}^\top.
\end{equation}

The effective weight change after the SGD step is $\Delta \mathbf{W}_{\text{LoRA}} = s(\mathbf{B}'\mathbf{A}' - \mathbf{B}\mathbf{A})$. Expanding the product:
\begin{align}
    \Delta \mathbf{W}_{\text{LoRA}}
    &= s\big[(\mathbf{B} - \eta s \mathbf{G}^\top \mathbf{X}\mathbf{A}^\top)(\mathbf{A} - \eta s \mathbf{B}^\top \mathbf{G}^\top \mathbf{X}) - \mathbf{B}\mathbf{A}\big] \nonumber\\
    &= s\big[\mathbf{B}\mathbf{A} - \eta s \mathbf{B}\mathbf{B}^\top \mathbf{G}^\top \mathbf{X} - \eta s \mathbf{G}^\top \mathbf{X}\mathbf{A}^\top \mathbf{A} + \eta^2 s^2 \mathbf{G}^\top \mathbf{X}\mathbf{A}^\top \mathbf{B}^\top \mathbf{G}^\top \mathbf{X} - \mathbf{B}\mathbf{A}\big] \nonumber\\
    &= -\eta s^2 \big(\mathbf{B}\mathbf{B}^\top \mathbf{G}^\top \mathbf{X} + \mathbf{G}^\top \mathbf{X}\mathbf{A}^\top \mathbf{A}\big) + \eta^2 s^3 \big(\mathbf{G}^\top \mathbf{X}\mathbf{A}^\top \mathbf{B}^\top \mathbf{G}^\top \mathbf{X}\big).
\end{align}

The second term is $O(\eta^2)$ since it scales quadratically with the learning rate. Under the standard assumption that $\eta \ll 1$ (small learning rate regime commonly used in fine-tuning), the quadratic term is negligible compared to the linear term. Dropping the $O(\eta^2)$ term yields Eq.~\eqref{eq:lora_update}.
\end{proof}

%==============================================================================
\section{PyTorch Pseudocode for LoRA-NSP}
\label{app:pseudocode}

We provide PyTorch-style pseudocode for the key components of LoRA-NSP in Algorithm~\ref{alg:lora_nsp}. The implementation comprises three stages per task: (i) online covariance estimation via forward hooks, (ii) soft projection matrix construction via eigendecomposition, and (iii) task training with the FD+CD distillation loss. After each task, LoRA weights are merged into the frozen backbone and reinitialized for the next task.

\begin{algorithm}[h]
\caption{LoRA-NSP Training Procedure (PyTorch-style)}
\label{alg:lora_nsp}
\begin{algorithmic}[1]
\Require Pre-trained CLIP model $f_I, f_T$; Task sequence $\{\mathcal{T}^1, \ldots, \mathcal{T}^K\}$; LoRA rank $r$; EMA momentum $\alpha$; NSP weight $\lambda_{\text{nsp}}$; Distillation coefficients $\lambda_{\text{fd}}, \lambda_{\text{cd}}$
\Ensure Fine-tuned image encoder with merged LoRA weights

\State \textbf{Initialize:} $W_0 \leftarrow \text{frozen pre-trained weights}$; $\Sigma_{\text{cov}}^{(l)} \leftarrow \mathbf{0}$ for each LoRA layer $l$

\For{each task $k = 1, \ldots, K$}
    \State \textcolor{gray}{\# Phase 1: Collect covariance statistics}
    \For{each batch $(\mathbf{x}, y)$ in $\mathcal{D}^k$}
        \State $\mathbf{X}_l^{\text{raw}} \leftarrow \text{extract\_input\_features}(\mathbf{x}, \text{layer}=l)$ \textcolor{gray}{\# via forward hook}
        \State $\mathbf{X}_l \leftarrow \text{flatten}(\mathbf{X}_l^{\text{raw}})$ \textcolor{gray}{\# $[n \cdot d_{\text{token}}, d_{\text{in}}]$}
        \State $\Sigma_{\text{cov}}^{(l)} \leftarrow \alpha \, \Sigma_{\text{cov}}^{(l)} + (1-\alpha) \, \mathbf{X}_l^\top \mathbf{X}_l$
    \EndFor

    \State \textcolor{gray}{\# Phase 2: Construct soft projection matrix}
    \For{each LoRA layer $l$}
        \State $\mathbf{U}, \boldsymbol{\Lambda} \leftarrow \text{eigendecomposition}(\Sigma_{\text{cov}}^{(l)})$
        \State $\mathbf{P}_{\text{SGP}}^{(l)} \leftarrow \sum_{i=1}^{d_{\text{in}}} w(\lambda_i) \, \mathbf{u}_i \mathbf{u}_i^\top$ \textcolor{gray}{\# Eq.~(5)}
    \EndFor

    \State \textcolor{gray}{\# Phase 3: Training with LoRA-NSP}
    \State Initialize $\mathbf{A}^{(l)} \sim \mathcal{N}(\mathbf{0}, \mathbf{I})$, $\mathbf{B}^{(l)} \leftarrow \mathbf{0}$ for each layer $l$
    \For{each epoch}
        \For{each batch $(\mathbf{x}, y)$ in $\mathcal{D}^k$}
            \State \textcolor{gray}{\# Forward pass with soft projection}
            \State $\mathbf{z} \leftarrow f_I(\mathbf{x})$ \textcolor{gray}{\# $W = W_0 + s \cdot \mathbf{B} \mathbf{A} \mathbf{P}_{\text{SGP}}$ per layer}
            \State $\text{logits} \leftarrow \mathbf{z} \cdot \mathbf{W}_{\text{text}}^\top$ \textcolor{gray}{\# zero-shot text embeddings}
            \State $\mathcal{L}_{\text{CE}} \leftarrow \text{CrossEntropy}(\text{logits}, y)$
            
            \State \textcolor{gray}{\# FD+CD distillation on reference data}
            \State $\mathcal{L}_{\text{fd}} \leftarrow \|f_I^{\text{pre}}(\mathbf{x}_{\text{ref}}) - f_I(\mathbf{x}_{\text{ref}})\|_2^2$
            \State $\mathcal{L}_{\text{cd}} \leftarrow \text{ContrastiveLoss}\big(f_I(\mathbf{x}_{\text{ref}}), f_T(\mathbf{t}_{\text{ref}})\big)$
            
            \State $\mathcal{L} \leftarrow \mathcal{L}_{\text{CE}} + \lambda_{\text{fd}} \mathcal{L}_{\text{fd}} + \lambda_{\text{cd}} \mathcal{L}_{\text{cd}}$
            \State $\mathcal{L}.\text{backward}()$; $\text{optimizer}.\text{step}()$
        \EndFor
    \EndFor

    \State \textcolor{gray}{\# Phase 4: Weight merging and reinitialization}
    \For{each LoRA layer $l$}
        \State $W_0^{(l)} \leftarrow W_0^{(l)} + s \cdot \mathbf{B}^{(l)} \mathbf{A}^{(l)} \mathbf{P}_{\text{SGP}}^{(l)}$
        \State Reset $\mathbf{A}^{(l)} \sim \mathcal{N}(\mathbf{0}, \mathbf{I})$, $\mathbf{B}^{(l)} \leftarrow \mathbf{0}$
    \EndFor
\EndFor

\State \Return $f_I$ with merged weights $W_0$
\end{algorithmic}
\end{algorithm}

\vspace{0.5em}
\noindent\textbf{Remarks on implementation.}
\begin{itemize}
    \item \textbf{Covariance extraction:} Input features to each LoRA layer are captured via \texttt{register\_forward\_hook} on the parent module. The raw tensor of shape $[n, d_{\text{token}}, d_{\text{in}}]$ is flattened along the batch and token dimensions to $[n \cdot d_{\text{token}}, d_{\text{in}}]$ before computing the Gram matrix.
    \item \textbf{Eigendecomposition frequency:} In practice, $\mathbf{P}_{\text{SGP}}$ is recomputed once per epoch (or every few epochs) rather than every batch, as the covariance evolves slowly.
    \item \textbf{Reference data for distillation:} We use Flickr8K as the reference dataset for FD+CD. The pre-trained model $f_I^{\text{pre}}$ is kept frozen in evaluation mode to provide stable target features.
    \item \textbf{Memory overhead:} The covariance matrix $\Sigma_{\text{cov}}^{(l)} \in \mathbb{R}^{d_{\text{in}} \times d_{\text{in}}}$ is the primary memory overhead. For CLIP ViT-B/16 with $d_{\text{in}} = 768$ and 12 transformer layers, this totals $12 \times 768^2 \times 4\,\text{bytes} \approx 28\,\text{MB}$ in float32, which is negligible compared to the model size.
\end{itemize}

\fi
\end{document}
